\section{AIME II 2021}
        \begin{enumerate}
        	\item (\textit{AIME II 2021}) Find the arithmetic mean of all the three-digit palindromes. (Recall that a palindrome is a number that reads the same forward and backward, such as $777$ or $383$.)


 
 \textbf{Answer: }550



\textbf{Solution 1}

Recall that the arithmetic mean of all the $n$ digit palindromes is just the average of the largest and smallest $n$ digit palindromes, and in this case the $2$ palindromes are $101$ and $999$ and $\frac{101+999}{2}=\boxed{550},$ which is the final answer.

 





\textbf{Solution 2}

For any palindrome $\underline{ABA},$ note that $\underline{ABA}$ is $100A + 10B + A = 101A + 10B.$ The average for $A$ is $5$ since $A$ can be any of $1, 2, 3, 4, 5, 6, 7, 8,$ or $9.$ The average for $B$ is $4.5$ since $B$ is either $0, 1, 2, 3, 4, 5, 6, 7, 8,$ or $9.$ Therefore, the answer is $505 + 45 = \boxed{550}.$

 - ARCTICTURN

 





\textbf{Solution 3 (Symmetry and Generalization)}

For every three-digit palindrome $\underline{ABA}$ with $A\in\{1,2,3,4,5,6,7,8,9\}$ and $B\in\{0,1,2,3,4,5,6,7,8,9\},$ note that $\underline{(10-A)(9-B)(10-A)}$ must be another palindrome by symmetry. Therefore, we can pair each three-digit palindrome uniquely with another three-digit palindrome so that they sum to 
 \begin{align*} \underline{ABA}+\underline{(10-A)(9-B)(10-A)}&=\left[100A+10B+A\right]+\left[100(10-A)+10(9-B)+(10-A)\right] \\ &=\left[100A+10B+A\right]+\left[1000-100A+90-10B+10-A\right] \\ &=1000+90+10 \\ &=1100. \end{align*}
For instances:
 \begin{align*} 171+929&=1100, \\ 262+838&=1100, \\ 303+797&=1100, \\ 414+686&=1100, \\ 545+555&=1100, \end{align*}
and so on.

 From this symmetry, the arithmetic mean of all the three-digit palindromes is $\frac{1100}{2}=\boxed{550}.$

 <u>\textbf{Remark}</u>

 By the Multiplication Principle, there are $9\cdot10=90$ three-digit palindromes in total. Their sum is $1100\cdot45=49500,$ as we can match them into $45$ pairs such that each pair sums to $1100.$

 





\textbf{Solution 4 (Similar to Solution 2: Very, Very Easy and Quick)}

We notice that a three-digit palindrome looks like this: $\underline{aba}.$

 And we know $a$ can be any digit from $1$ through $9,$ and $b$ can be any digit from $0$ through $9,$ so there are $9\times{10}=90$ three-digit palindromes.

 We want to find the sum of these $90$ palindromes and divide it by $90$ to find the arithmetic mean.

 How can we do that? Instead of adding the numbers up, we can break each palindrome into two parts: $101a+10b.$

 Thus, all of these $90$ palindromes can be broken into this form.

 Thus, the sum of these $90$ palindromes will be $101\times{(1+2+...+9)}\times{10}+10\times{(0+1+2+...+9)}\times{9},$ because each $a$ will be in $10$ different palindromes (since for each $a,$ there are $10$ choices for $b$). The same logic explains why we multiply by $9$ when computing the total sum of $b.$

 We get a sum of $45\times{1100},$ but don't compute this! Divide this by $90$ and you will get $\boxed{550}.$

 





\textbf{Solution 5 (Extremely Fast Solution)}

The possible values of the first and last digits each are $1, 2, ..., 8, 9$ with a sum of $45$ so the average value is $5.$ The middle digit can be any digit from $0$ to $9$ with a sum of $45,$ so the average value is $4.5.$ The average of all three-digit palindromes is $5\cdot 10^2+4.5\cdot 10+5=\boxed{550}.$

 





\textbf{Solution 6 (Two cases)}

<i>\textbf{Case 1}</i>

 Consider palindromes of the form $5x5 = 505 + 10x.$ There are $10$ of them. The arithmetic mean of the first term is $505,$ and the second $\frac {10 \cdot(0 + 1 + ... + 9)}{10} = 45.$The arithmetic mean of the sum is $505 + 45 = 550.$

 <i>\textbf{Case 2}</i>

 Consider palindromes of the form $yxy,$ where $y= {1,2,3,4,6,7,8,9}.$Let $u = 10 – y, v = 9 – x.$Then $uvu$ is a symmetric palindrome that can be constructed for each $yxy.$The arithmetic mean of each such pair is $550.$ For example, $\frac{737 + 363}{2} = 550.$

 Thus, all palindromes are divided into groups of numbers with the arithmetic mean in each group equal to $550.$ 

 The arithmetic mean of all numbers is also $550.$

 \textbf{vladimir.shelomovskii@gmail.com, vvsss}

 \textbf{Remark} Visit the \href{https://artofproblemsolving.com/wiki/index.php/Talk:2021\_AIME\_II\_Problems/Problem\_1}{Discussion Page} for questions and further generalizations.

 speedy 2 min video

 




	\item (\textit{AIME II 2021}) Equilateral triangle $ABC$ has side length $840$. Point $D$ lies on the same side of line $BC$ as $A$ such that $\overline{BD} \perp \overline{BC}$. The line $\ell$ through $D$ parallel to line $BC$ intersects sides $\overline{AB}$ and $\overline{AC}$ at points $E$ and $F$, respectively. Point $G$ lies on $\ell$ such that $F$ is between $E$ and $G$, $\triangle AFG$ is isosceles, and the ratio of the area of $\triangle AFG$ to the area of $\triangle BED$ is $8:9$. Find $AF$.

                    \begin{figure}[H]
                        \centering
                        \includegraphics[width=0.5\textwidth]{Problem/AIME_2021_II_2_Problem_Figure_1.png}
                    \end{figure}
                    


 
 \textbf{Answer: }336



\textbf{Solution 1 (Area Formulas for Triangles)}

By angle chasing, we conclude that $\triangle AGF$ is a $30^\circ\text{-}30^\circ\text{-}120^\circ$ triangle, and $\triangle BED$ is a $30^\circ\text{-}60^\circ\text{-}90^\circ$ triangle.

 Let $AF=x.$ It follows that $FG=x$ and $EB=FC=840-x.$ By the side-length ratios in $\triangle BED,$ we have $DE=\frac{840-x}{2}$ and $DB=\frac{840-x}{2}\cdot\sqrt3.$

 Let the brackets denote areas. We have 
 \[[AFG]=\frac12\cdot AF\cdot FG\cdot\sin{\angle AFG}=\frac12\cdot x\cdot x\cdot\sin{120^\circ}=\frac12\cdot x^2\cdot\frac{\sqrt3}{2}\] and 
 \[[BED]=\frac12\cdot DE\cdot DB=\frac12\cdot\frac{840-x}{2}\cdot\left(\frac{840-x}{2}\cdot\sqrt3\right).\]

 We set up and solve an equation for $x:$
 \begin{align*} \frac{[AFG]}{[BED]}&=\frac89 \\ \frac{\frac12\cdot x^2\cdot\frac{\sqrt3}{2}}{\frac12\cdot\frac{840-x}{2}\cdot\left(\frac{840-x}{2}\cdot\sqrt3\right)}&=\frac89 \\ \frac{2x^2}{(840-x)^2}&=\frac89 \\ \frac{x^2}{(840-x)^2}&=\frac49. \end{align*}
Since $0<x<840,$ it is clear that $\frac{x}{840-x}>0.$ Therefore, we take the positive square root for both sides:
 \begin{align*} \frac{x}{840-x}&=\frac23 \\ 3x&=1680-2x \\ 5x&=1680 \\ x&=\boxed{336}. \end{align*}


 





\textbf{Solution 2}

We express the areas of $\triangle BED$ and $\triangle AFG$ in terms of $AF$ in order to solve for $AF.$

 We let $x = AF.$ Because $\triangle AFG$ is isosceles and $\triangle AEF$ is equilateral, $AF = FG = EF = AE = x.$ 

 Let the height of $\triangle ABC$ be $h$ and the height of $\triangle AEF$ be $h'.$ Then we have that $h = \frac{\sqrt{3}}{2}(840) = 420\sqrt{3}$ and $h' = \frac{\sqrt{3}}{2}(EF) = \frac{\sqrt{3}}{2}x.$ 

 Now we can find $DB$ and $BE$ in terms of $x.$ $DB = h - h' = 420\sqrt{3} - \frac{\sqrt{3}}{2}x,$ $BE = AB - AE = 840 - x.$ Because we are given that $\angle DBC = 90,$ $\angle DBE = 30.$ This allows us to use the sin formula for triangle area: the area of $\triangle BED$ is $\frac{1}{2}(\sin 30)\left(420\sqrt{3} - \frac{\sqrt{3}}{2}x\right)(840-x).$ Similarly, because $\angle AFG = 120,$ the area of $\triangle AFG$ is $\frac{1}{2}(\sin 120)(x^2).$

 Now we can make an equation:
 \begin{align*} \frac{\triangle AFG}{\triangle BED} &= \frac{8}{9} \\ \frac{\frac{1}{2}(\sin 120)(x^2)}{\frac{1}{2}(\sin 30)\left(420\sqrt{3} - \frac{\sqrt{3}}{2}x\right)(840-x)} &= \frac{8}{9} \\ \frac{x^2}{\left(420 - \frac{x}{2}\right)(840-x)} &= \frac{8}{9}. \end{align*}
To make further calculations easier, we scale everything down by $420$ (while keeping the same variable names, so keep that in mind). 
 \begin{align*} \frac{x^2}{\left(1-\frac{x}{2}\right)(2-x)} &= \frac{8}{9} \\ 8\left(1-\frac{x}{2}\right)(2-x) &= 9x^2 \\ 16-16x + 4x^2 &= 9x^2 \\ 5x^2 + 16x -16 &= 0 \\ (5x-4)(x+4) &= 0. \end{align*}
Thus $x = \frac{4}{5}.$ Because we scaled down everything by $420,$ the actual value of $AF$ is $\frac{4}{5}(420) = \boxed{336}.$

 





\textbf{Solution 3 (Pretty Straightforward)}

$\angle AFE = \angle AEF = \angle EAF = 60^{0} \Rightarrow \angle AFG = 120^{0}$So, If $\Delta AFG$ is isosceles, it means that $AF = FG$.

 Let $AF = FG = AE = EF = x$

 So, $[\Delta AFG] = \frac{1}{2} \cdot x^{2} \textup{sin} 120^{0} = \frac{\sqrt{3}}{4}x^{2}$

 In $\Delta BED$, $BE = 840 - x$, Hence $DE = \frac{840 - x}{2}$ (because $\textup{sin} 30^{0} = \frac{1}{2}$)

 Therefore, $[\Delta BED] = \frac{1}{2} (840 - x) \left (\frac{840-x}{2} \right) \textup{sin} 60^{0}$

 So, $[\Delta BED] = \frac{\sqrt{3}}{4} (840 - x) \left (\frac{840-x}{2} \right) = \frac{\sqrt{3}}{8} (840 - x)^{2}$

 
Now, as we know that the ratio of the areas of $\Delta AFG$ and $\Delta BED$ is $8:9$

 Substituting the values, we get

 $\frac{\frac{\sqrt{3}}{4}x^{2}}{\frac{\sqrt{3}}{8} (840 - x)^{2}} = \frac{8}{9} \Rightarrow \left (\frac{x}{840 - x} \right)^{2} = \frac{4}{9}$Hence, $\frac{x}{840 - x} = \frac{2}{3}$. Solving this, we easily get $x = 336$

 We have taken $AF = x$, Hence, $AF = \boxed{336}$

 -Arnav Nigam

 





\textbf{Solution 4 (Similar Triangles)}

Since $\triangle AFG$ is isosceles, $AF = FG$, and since $\triangle AEF$ is equilateral, $AF = EF$. Thus, $EF = FG$, and since these triangles share an altitude, they must have the same area.

 Drop perpendiculars from $E$ and $F$ to line $BC$; call the meeting points $P$ and $Q$, respectively. $\triangle BEP$ is clearly congruent to both $\triangle BED$ and $\triangle FQC$, and thus each of these new triangles has the same area as $\triangle BED$. But we can "slide" $\triangle BEP$ over to make it adjacent to $\triangle FQC$, thus creating an equilateral triangle whose area has a ratio of $18:8$ when compared to $\triangle AEF$ (based on our conclusion from the first paragraph). Since these triangles are both equilateral, they are similar, and since the area ratio $18:8$ reduces to $9:4$, the ratio of their sides must be $3:2$. So, because $FC$ and $AF$ represent sides of these triangles, and they add to $840$, $AF$ must equal two-fifths of $840$, or $\boxed{336}$.

 





\textbf{Solution 5 (No Trig, Straightforward)}

Draw the perpendicular segment from $A$ to $BC$. Call its intersection with $EF$ as point $H$ and its intersection with $BC$ as point $I$. Because $AFH$ is a $30-60-90$ triangle, $AH = \frac{\sqrt{3}}{2}x$. Call $DB = HI$ as $h$. Because $ABI$ is a $30-60-90$ triangle, $AI = \frac{\sqrt{3}}{2}x + h = 420\sqrt{3}$. $DE = \frac{h}{\sqrt{3}}$ because $DBE$ is a $30-60-90$ triangle. We know from the problem statement that $9[AFG]= \frac{9}{2}(\frac{\sqrt{3}}{2}x^2) = \frac{8}{2\sqrt{3}}h^2 = 8[BED]$. By simplifying this equation, $h = \frac{3\sqrt{3}x}{4}$. Substituting $h$ into the first equation yields $x = \boxed{336}$.

 -unhappyfarmer

 





\textbf{Solution 6 (One Altitude & No Radicals)}

Drop an altitude from $F$ to $\overline{AG}$ which intersects $\overline{AG}$ at point $H$. Both triangles $\triangle DEB$ and $\triangle FHA$ are $30$-$60$-$90$ triangles, and therefore similar.

 Since altitude $\overline{FH}$ bisects the area of $\triangle AFG$, the area of $\triangle FHA$ is exactly $\frac{1}{2}$ of the area of $\triangle AFG$. The ratio of the area of $\triangle FHA$ to $\triangle DEB$ is then $\frac{4}{9}$, so the ratio of their corresponding sides is $\frac{2}{3}$. 

 Since triangle $\triangle AEF$ is equilateral, $\overline{AE} = \overline{AF} = \frac{2}{3}  \overline{EB} \implies \overline{AE} + \overline{EB} = \frac{5}{3} \overline{EB} = 840$. 

 Therefore, $\overline{EB} = 540$, and $\overline{AB} = \boxed{336}$ 

 





\textbf{Solution 7}

Drop an altitude from $F$ to $BC$ and notice that the triangle is congruent to triangle $DBE$, so this triangle also has a ratio of $\frac{9}{8}$ with respect to triangle $AFG$. We can let $AF$ and $FG$ be x and let $FC$ be $(840-x)$. We can split $AFG$ into 2 $30$-$60$-$90$ triangles with hypotenuse $x$. The base of these triangles is $\frac{x}{2}$, and since each of these triangles is similar to the triangle we just crafted (let this triangle be $FCZ$), we have a neat way to find the value of $x$. We can say that the base of one of the new $30$-$60$-$90$ triangles we crafted is $\frac{x}{2}$, and since the hypotenuse is $x$, we don't need to work with any radicals at all! We can multiply the hypotenuse by the base in each of the smaller triangles of $AFG$, add the values together, put the result as the numerator of the fraction with the hypotenuse times the base of $FCZ$, and set this expression equal to $\frac{8}{9}$.

 When we actually work this out, we get that the numerator will be $(x^2)$ and the denominator will be $\frac{(840-x)^2}{2}$, and we set this equal to $\frac{8}{9}$.

 We then solve the equation for x.

 \begin{align*}\frac{2x^2}{(840-x)^2} &= \frac{8}{9} \\[10pt]\frac{x^2}{(840-x)^2} &= \frac{4}{9} \\[10pt]\sqrt{\frac{x^2}{(840-x)^2}} &= \sqrt{\frac{4}{9}} \\[10pt]\left| \frac{x}{840-x} \right| &= \frac{2}{3} \\[10pt]\frac{x}{840-x} &= \pm \frac{2}{3}\end{align*}


 \begin{array}{rcclr}\frac{x}{840-x} &= \frac{2}{3} &\quad \lor \quad& \frac{x}{840-x} &= -\frac{2}{3} \\[10pt]3x &= 2(840-x) &\quad \lor \quad& 3x &= -2(840-x) \\[10pt]3x &= 1680 - 2x &\quad \lor \quad& 3x &= -1680 + 2x \\[10pt]5x &= 1680 &\quad \lor \quad& x &= -1680 \\[10pt]x &= 336 &\quad \lor \quad& x &= -1680\end{array}

 Clearly the only value of x that works is $\boxed{336}$

 




	\item (\textit{AIME II 2021}) Find the number of permutations $x_1, x_2, x_3, x_4, x_5$ of numbers $1, 2, 3, 4, 5$ such that the sum of five products 
 \[x_1x_2x_3 + x_2x_3x_4 + x_3x_4x_5 + x_4x_5x_1 + x_5x_1x_2\] is divisible by $3$.


 
 \textbf{Answer: }080



\textbf{Solution 1}

Since $3$ is one of the numbers, a product with a $3$ in it is automatically divisible by $3,$ so WLOG $x_3=3,$ and we will multiply by $5$ afterward since any of $x_1, x_2, \ldots, x_5$ would be $3,$ after some cancelation we see that now all we need to find is the number of ways that $x_5x_1(x_4+x_2)$ is divisible by $3,$ since $x_5x_1$ is never divisible by $3,$ now we just need to find the number of ways $x_4+x_2$ is divisible by $3.$ Note that $x_2$ and $x_4$ can be $(1, 2), (2, 1), (1, 5), (5, 1), (2, 4), (4, 2), (4, 5),$ or $(5, 4).$ We have $2$ ways to designate $x_1$ and $x_5$ for a total of $8 \cdot 2 = 16.$ So the desired answer is $16 \cdot 5=\boxed{080}.$

 





\textbf{Solution 2 (Symmetry and Casework)}

The expression $x_1x_2x_3 + x_2x_3x_4 + x_3x_4x_5 + x_4x_5x_1 + x_5x_1x_2$ has symmetry. Without the loss of generality, let $x_1=3.$ It follows that $\{x_2,x_3,x_4,x_5\}=\{1,2,4,5\}.$ We have:

 <ol style="margin-left: 1.5em;"> \item $x_1x_2x_3 + x_2x_3x_4 + x_3x_4x_5 + x_4x_5x_1 + x_5x_1x_2\equiv x_2x_3x_4 + x_3x_4x_5\pmod{3}.$ 

 \item $x_2,x_3,x_4,x_5$ are congruent to $1,2,1,2\pmod{3}$ in some order.

 </ol>We construct the following table for the case $x_1=3,$ with all values in modulo $3:$
 \[\begin{array}{c||c|c|c|c|c||c} & & & & & & \\ [-2.5ex] \textbf{Row} & \boldsymbol{x_2} & \boldsymbol{x_3} & \boldsymbol{x_4} & \boldsymbol{x_5} & \boldsymbol{x_2x_3x_4 + x_3x_4x_5} & \textbf{Valid?} \\ [0.5ex] \hline & & & & & & \\ [-2ex] 1 & 1 & 1 & 2 & 2 & 0 & \checkmark \\ 2 & 1 & 2 & 1 & 2 & 0 & \checkmark \\ 3 & 1 & 2 & 2 & 1 & 2 & \\ 4 & 2 & 1 & 1 & 2 & 1 & \\ 5 & 2 & 1 & 2 & 1 & 0 & \checkmark \\ 6 & 2 & 2 & 1 & 1 & 0 & \checkmark \end{array}\]For Row 1, $(x_2,x_3)$ can be either $(1,4)$ or $(4,1),$ and $(x_4,x_5)$ can be either $(2,5)$ or $(5,2).$ By the Multiplication Principle, Row 1 produces $2\cdot2=4$ permutations. Similarly, Rows 2, 5, and 6 each produce $4$ permutations.

 Together, we get $4\cdot4=16$ permutations for the case $x_1=3.$ By symmetry, the cases $x_2=3, x_3=3, x_4=3,$ and $x_5=3$ all have the same count. Therefore, the total number of permutations $x_1, x_2, x_3, x_4, x_5$ is $16\cdot5=\boxed{080}.$

 





\textbf{Solution 3}

WLOG, let $x_{3} = 3$So, the terms $x_{1}x_{2}x_{3}, x_{2}x_{3}x_{4},x_{3}x_{4}x_{5}$ are divisible by $3$.

 We are left with $x_{4}x_{5}x_{1}$ and $x_{5}x_{1}x_{2}$.We need $x_{4}x_{5}x_{1} + x_{5}x_{1}x_{2} \equiv 0 \pmod{3}$.The only way is when They are $(+1,-1)$ or $(-1, +1) \pmod{3}$.

 The numbers left with us are $1,2,4,5$ which are $+1,-1,+1,-1\pmod{3}$ respectively.

 $+1$ (of $x_{4}x_{5}x_{1}$ or $x_{5}x_{1}x_{2}$) $= +1 \cdot +1 \cdot +1$ $\;\;\; OR \;\;\;+1$ (of $x_{4}x_{5}x_{1}$ or $x_{5}x_{1}x_{2}$) = $-1 \cdot -1 \cdot +1$.

 $-1$ (of $x_{4}x_{5}x_{1}$ or $x_{5}x_{1}x_{2}$) $= -1 \cdot -1 \cdot -1$ $\;\;\; OR \;\;\;-1$ (of $x_{4}x_{5}x_{1}$ or $x_{5}x_{1}x_{2}$) = $-1 \cdot +1 \cdot +1$

 But, as we have just two $+1's$ and two $-1's$.Hence, We will have to take $+1 = +1 \cdot -1 \cdot -1$ and $-1 = -1 \cdot +1 \cdot +1$.Among these two, we have a $+1$ and $-1$ in common, i.e. $(x_{5}, x_{1}) = (+1, -1) or (-1, +1)$ (because $x_{1}$ and $x_{5}$.are common in $x_{4}x_{5}x_{1}$ and $x_{5}x_{1}x_{2}$).

 So, $(x_{5}, x_{1}) \in {(1,2), (1,5), (4,2), (4,5), (2,1), (5,1), (2,4), (5,4)}$ i.e. $8$ values.

 For each value of $(x_{5}, x_{1})$ we get $2$ values for $(x_{2}, x_{4})$.Hence, in total, we have $8 \times 2 = 16$ ways.

 But any of the $x_{i} 's$ can be $3$.So, $16 \times 5 = \boxed{080}$.

 -Arnav Nigam

 





\textbf{Solution 4 (Proportion)}

WLOG, let $x_{3} = 3$. Then:
 \[x_{1}x_{2}x_{3} + x_{2}x_{3}x_{4} + x_{3}x_{4}x_{5} + x_{4}x_{5}x_{1} + x_{5}x_{1}x_{2} = 3 (x_1 x_2 + x_2 x_4 + x_4 x_5) + x_5 x_1 (x_2  + x_4).\]The sum is divisible by $3$ if and only if  $x_2  + x_4$  is divisible by $3$.The possible sums of  $x_2  + x_4$ are $1 + 2, 1 + 4, 1 + 5, 2 + 4, 2 + 5, 4 + 5.$Two of them are not multiples of $3$, but four of them are multiples.

 A total number of  permutations is $5! = 120.$ 

 $\frac {2}{3}$ of this number, that is, $80,$ give sums that are multiples of $3.$

 \textbf{vladimir.shelomovskii@gmail.com, vvsss}

 





\textbf{Solution 5 (Factoring)}

This is my first time doing a solution (feel free to edit it)

 We have 
 \[x_1x_2x_3 + x_2x_3x_4 + x_3x_4x_5 + x_4x_5x_1 + x_5x_1x_2.\]

 We have $5$ numbers. Considering any $x$ as $3,$ we see that we are left with two terms that are not always divisible by $3,$ which means that already gives us 5 options. Let's now consider $x_1 =3:$We are left with 
 \[3x_2x_3 + x_2x_3x_4 + x_3x_4x_5 + 3x_4x_5 + 3x_5x_2.\]The two terms left over are 
 \[x_2x_3x_4 + x_3x_4x_5 \equiv 0 \pmod{3}\]since we already have used $3$ the remaining numbers are $1,2,4,5.$

 We now factor 
 \begin{align*} (x_2 + x_5)(x_3x_4) &\equiv 0 \pmod{3} \\ (x_2 + x_5) &\equiv 0 \pmod{3} \end{align*}
since $1,2,4,5$ are all not factors of $3.$

 Now using the number $1,2,4,5,$ we take two to get a number divisible by $3$ for $(x_2 + x_5):$
 \begin{align*} 1+5 &\equiv 0 \pmod{3}, \\ 4+2 &\equiv 0 \pmod{3}, \\ 4+5 &\equiv 0 \pmod{3}, \\ 1+2 &\equiv 0 \pmod{3}. \end{align*}
We have $4$ possibilities from above. Since we can also have $5+1$ or $2+4,$ there are $4\cdot2=8$ possibilities in all.

 Now using 
 \[(x_2 + x_5)(x_3x_4) \equiv 0 \pmod{3},\] we have $(x_3x_4),$ which results in $8$ more possibilities of $2$ times more. So, we get $2\cdot2\cdot4=16.$

 Remember that $3$ can be any of $5$ different variables. So, we multiply by $5$ to get the answer $\boxed{080}.$

 




	\item (\textit{AIME II 2021}) There are real numbers $a, b, c,$ and $d$ such that $-20$ is a root of $x^3 + ax + b$ and $-21$ is a root of $x^3 + cx^2 + d.$ These two polynomials share a complex root $m + \sqrt{n} \cdot i,$ where $m$ and $n$ are positive integers and $i = \sqrt{-1}.$ Find $m+n.$


 
 \textbf{Answer: }330



\textbf{Solution 1 (Complex Conjugate Root Theorem and Vieta's Formulas)}

By the Complex Conjugate Root Theorem, the imaginary roots for each of $x^3+ax+b$ and $x^3+cx^2+d$ are complex conjugates. Let $z=m+\sqrt{n}\cdot i$ and $\overline{z}=m-\sqrt{n}\cdot i.$ It follows that the roots of $x^3+ax+b$ are $-20,z,\overline{z},$ and the roots of $x^3+cx^2+d$ are $-21,z,\overline{z}.$

 We know that
 \begin{align*} z+\overline{z}&=2m, & (1) \\ z\overline{z}&=m^2+n. & (2) \end{align*}
Applying Vieta's Formulas to $x^3+ax+b,$ we have $-20+z+\overline{z}=0.$ Substituting $(1)$ into this equation, we get $m=10.$

 Applying Vieta's Formulas to $x^3+cx^2+d,$ we have $-21z-21\overline{z}+z\overline{z}=0,$ or $-21(z+\overline{z})+z\overline{z}=0.$ Substituting $(1)$ and $(2)$ into this equation, we get $n=320.$

 Finally, the answer is $m+n=\boxed{330}.$

 





\textbf{Solution 2 (Somewhat Bashy)}

$(-20)^{3} + (-20)a + b = 0$, hence $-20a + b = 8000$

 Also, $(-21)^{3} + c(-21)^{2} + d = 0$, hence $441c + d = 9261$

 $m + i \sqrt{n}$ satisfies both $\Rightarrow$ we can put it in both equations and equate to 0.

 In the first equation, we get$(m + i \sqrt{n})^{3} + a(m + i \sqrt{n}) + b = 0$Simplifying this further, we get $(m^{3} - 3mn + am + b) + i(3m^{2} \sqrt{n} - n\sqrt{n} + a\sqrt{n}) = 0$

 Hence, $m^{3} - 3mn + am + b = 0$ and $3m^{2} \sqrt{n} - n\sqrt{n} + a\sqrt{n} = 0 \Rightarrow 3m^{2} - n + a = 0 \rightarrow (1)$

 In the second equation, we get$(m + i \sqrt{n})^{3} + c(m + i \sqrt{n})^{2} + d = 0$Simplifying this further, we get $(m^{3} + m^{2}c - nc - 3mn + d) + i(3m^{2} \sqrt{n} - n\sqrt{n} + 2mc\sqrt{n}) = 0$

 Hence, $m^{3} + m^{2}c - nc - 3mn + d = 0$ and $3m^{2} \sqrt{n} - n\sqrt{n} + 2mc\sqrt{n} = 0 \Rightarrow 3m^{2} - n + 2mc = 0 \rightarrow (2)$

 Comparing (1) and (2),

 $a = 2mc$ and $am + b = m^{2}c - nc + d \rightarrow (3)$

 $b = 8000 + 20a \Rightarrow b = 40mc + 8000$; $d = 9261 - 441c$

 Substituting these in $(3)$ gives, $2m^{2}c + 8000 + 40mc = m^{2}c - nc + 9261 - 441c$

 This simplifies to $m^{2}c + nc + 40mc + 441c = 1261 \Rightarrow c(m^{2} + n + 40m + 441) = 1261$

 Hence, $c|1261 \Rightarrow c \in {1,13,97,1261}$

 
Consider case of $c = 1$:

 $c = 1 \Rightarrow d = 8820$Also, $a = 2m, b = 8000 + 40m$

 $am + b = m^{2} - n + 8820$ (because c = 1)Also, $m^{2} + n + 40m = 820 \rightarrow (4)$Also, Equation (2) gives $3m^{2} - n + 2m = 0 \rightarrow (5)$

 Solving (4) and (5) simultaneously gives $m = 10, n = 320$

 [AIME can not have more than one answer, so we can stop here ... Not suitable for objective exam]

 Hence, $m + n = 10 + 320 = \boxed{330}$

 -Arnav Nigam

 





\textbf{Solution 3 (Heavy Calculation Solution)}

We start off by applying Vieta's, and we find that $a=m^2+n-40m$, $b=20m^2+20n$, $c=21-2m$, and $d=21m^2+21n$. After that, we use the fact that $-20$ and $-21$ are roots of $x^3+ax+b$ and $x^3+cx^2+d$, respectively. Since substituting the roots back into the function returns zero, we have that $(-20)^3-20a+b=0$ and $(-21)^3+c\cdot (-21)^2+d=0$. Setting these two equations equal to each other while also substituting the values of $a$, $b$, $c$, and $d$ above gives us $21m^2+21n-1682m+8000=0$. We then rearrange the equation into $21n = -21m^2+1682m-8000$.

 With this property, we know that $-21m^2+1682m-8000$ is divisible by $21$, so $1682m-8000=0 \pmod{21}$. This results in $2m-20=0 \pmod{21}$, which finally gives us $m=10 \pmod{21}$. We can test the first obvious value of $m$, which is $10$, to find that this works, and we get $m=10$ and $n=320$. Therefore, the answer is $m + n = 10 + 320 = \boxed{330}.$

 





\textbf{Solution 4 (Synthetic Division)}

We note that $x^3 + ax + b = (x+20)P(x)$ and $x^3 + cx^2 + d = (x+21)Q(x)$ for some polynomials $P(x)$ and $Q(x)$.

 Through synthetic division (ignoring the remainder as we can set $b$ and $d$ to constant values such that the remainder is zero), $P(x) = x^2 - 20x + (400+a)$, and $Q(x) = x^2 + (c-21)x + (441 - 21c)$.

 By the complex conjugate root theorem, we know that $P(x)$ and $Q(x)$ share the same roots, and they share the same leading coefficient, so $P(x) = Q(x)$. 

 Therefore, $c-21 = -20$ and $441-21c = 400 + a$. Solving the system of equations, we get $a = 20$ and $c = 1$, so $P(x) = Q(x) = x^2 - 20x + 420$.

 Finally, by the quadratic formula, we have roots of $\frac{20 \pm \sqrt{400 - 1680}}{2} = 10 \pm \sqrt{320}i$, so our final answer is $10 + 320 = \boxed{330}$

 -faefeyfa

 





\textbf{Solution 5 (Fast and Easy)}

We plug -20 into the equation obtaining $(-20)^3-20a+b$, likewise, plugging -21 into the second equation gets $(-21)^3+441c+d$.

 





\textbf{Solution 6 (solution by integralarefun)}

Since $m+i\sqrt{n}$ is a common root and all the coefficients are real, $m-i\sqrt{n}$ must be a common root, too.

 Now that we know all three roots of both polynomials, we can match coefficients (or more specifically, the zero coefficients).

 First, however, the product of the two common roots is:
 \begin{align*} &&&(x-m-i\sqrt{n})(x-m+i\sqrt{n})\\ &=&&x^2-x(m+i\sqrt{n}+m-i\sqrt{n})+(m+i\sqrt{n})(m-i\sqrt{n})\\ &=&&x^2-2xm+(m^2-i^2n)\\ &=&&x^2-2xm+m^2+n \end{align*}


 Now, let's equate the two forms of both the polynomials:
 \[x^3+ax+b=(x^2-2xm+m^2+n)(x+20)\]
 \[x^3+cx^2+d=(x^2-2xm+m^2+n)(x+21)\]Now we can match the zero coefficients.
 \[-2m+20=0\to m=10\text{ and}\]
 \[-42m+m^2+n=0\to-420+100+n=0\to n=320\text{.}\]Thus, $m+n=10+320=\boxed{330}$.

 




	\item (\textit{AIME II 2021}) For positive real numbers $s$, let $\tau(s)$ denote the set of all obtuse triangles that have area $s$ and two sides with lengths $4$ and $10$. The set of all $s$ for which $\tau(s)$ is nonempty, but all triangles in $\tau(s)$ are congruent, is an interval $[a,b)$. Find $a^2+b^2$.


 
 \textbf{Answer: }736



\textbf{Solution 1}

We start by defining a triangle. The two small sides MUST add to a larger sum than the long side. We are given $4$ and $10$ as the sides, so we know that the $3$rd side is between $6$ and $14$, exclusive. We also have to consider the word OBTUSE triangles. That means that the two small sides squared is less than the $3$rd side. So the triangles' sides are between $6$ and $\sqrt{84}$ exclusive, and the larger bound is between $\sqrt{116}$ and $14$, exclusive. The area of these triangles are from $0$ (straight line) to $2\sqrt{84}$ on the first "small bound" and the larger bound is between $0$ and $20$. $0 < s < 2\sqrt{84}$ is our first equation, and $0 < s < 20$ is our $2$nd equation. Therefore, the area is between $\sqrt{336}$ and $\sqrt{400}$, so our final answer is $\boxed{736}$.

 





\textbf{Solution 2 (Inequalities and Casework)}

If $a,b,$ and $c$ are the side-lengths of an obtuse triangle with $a\leq b\leq c,$ then both of the following must be satisfied:

 
                \begin{itemize}
                    \item \textit{Triangle Inequality Theorem:} $a+b>c$
                \end{itemize}


                
                \begin{itemize}
                    \item \textit{Pythagorean Inequality Theorem:} $a^2+b^2<c^2$
                \end{itemize}


                For one such obtuse triangle, let $4,10,$ and $x$ be its side-lengths and $K$ be its area. We apply casework to its longest side:

 \textbf{Case (1): The longest side has length $\boldsymbol{10,}$ so $\boldsymbol{0<x<10.}$}

 By the Triangle Inequality Theorem, we have $4+x>10,$ from which $x>6.$

 By the Pythagorean Inequality Theorem, we have $4^2+x^2<10^2,$ from which $x<\sqrt{84}.$

 Taking the intersection produces $6<x<\sqrt{84}$ for this case.

 At $x=6,$ the obtuse triangle degenerates into a straight line with area $K=0;$ at $x=\sqrt{84},$ the obtuse triangle degenerates into a right triangle with area $K=\frac12\cdot4\cdot\sqrt{84}=2\sqrt{84}.$ Together, we obtain $0<K<2\sqrt{84},$ or $K\in\left(0,2\sqrt{84}\right).$

 \textbf{Case (2): The longest side has length $\boldsymbol{x,}$ so $\boldsymbol{x\geq10.}$}

 By the Triangle Inequality Theorem, we have $4+10>x,$ from which $x<14.$

 By the Pythagorean Inequality Theorem, we have $4^2+10^2<x^2,$ from which $x>\sqrt{116}.$

 Taking the intersection produces $\sqrt{116}<x<14$ for this case.

 At $x=14,$ the obtuse triangle degenerates into a straight line with area $K=0;$ at $x=\sqrt{116},$ the obtuse triangle degenerates into a right triangle with area $K=\frac12\cdot4\cdot10=20.$ Together, we obtain $0<K<20,$ or $K\in\left(0,20\right).$

 \textbf{Answer}

 It is possible for noncongruent obtuse triangles to have the same area. Therefore, all such positive real numbers $s$ are in exactly one of $\left(0,2\sqrt{84}\right)$ or $\left(0,20\right).$ By the exclusive disjunction, the set of all such $s$ is 
 \[[a,b)=\left(0,2\sqrt{84}\right)\oplus\left(0,20\right)=\left[2\sqrt{84},20\right),\] from which $a^2+b^2=\boxed{736}.$

 





\textbf{Solution 3}

We have the diagram below.

 
                    \begin{figure}[H]
                        \centering
                        \includegraphics[width=0.5\textwidth]{Solution/AIME_2021_II_5_Solution_Figure_1.png}
                    \end{figure}
                    

 We proceed by taking cases on the angles that can be obtuse, and finding the ranges for $s$ that they yield .

 If angle $\theta$ is obtuse, then we have that $s \in (0,20)$. This is because $s=20$ is attained at $\theta = 90^{\circ}$, and the area of the triangle is strictly decreasing as $\theta$ increases beyond $90^{\circ}$. This can be observed from
 \[s=\frac{1}{2}(4)(10)\sin\theta\]by noting that $\sin\theta$ is decreasing in $\theta \in (90^{\circ},180^{\circ})$.

 Then, we note that if $\alpha$ is obtuse, we have $s \in (0,4\sqrt{21})$. This is because we get $x=\sqrt{10^2-4^2}=\sqrt{84}=2\sqrt{21}$ when $\alpha=90^{\circ}$, yielding $s=4\sqrt{21}$. Then, $s$ is decreasing as $\alpha$ increases by the same argument as before.

 $\angle{ACB}$ cannot be obtuse since $AC>AB$.

 Now we have the intervals $s \in (0,20)$ and $s \in (0,4\sqrt{21})$ for the cases where $\theta$ and $\alpha$ are obtuse, respectively. We are looking for the $s$ that are in exactly one of these intervals, and because $4\sqrt{21}<20$, the desired range is
 \[s\in [4\sqrt{21},20)\]giving 
 \[a^2+b^2=\boxed{736}.         \Box\]

 





\textbf{Solution 4}

Note: Archimedes15 Solution which I added an answerhere are two cases. Either the $4$ and $10$ are around an obtuse angle or the $4$ and $10$ are around an acute triangle. If they are around the obtuse angle, the area of that triangle is $<20$ as we have $\frac{1}{2} \cdot 40 \cdot \sin{\alpha}$ and $\sin$ is at most $1$. Note that for the other case, the side lengths around the obtuse angle must be $4$ and $x$ where we have $16+x^2 < 100 \rightarrow x < 2\sqrt{21}$. Using the same logic as the other case, the area is at most $4\sqrt{21}$. Square and add $4\sqrt{21}$ and $20$ to get the right answer 
 \[a^2+b^2= \boxed{736}\Box\]

 





\textbf{Solution 5 (Circles)}

For $\triangle ABC,$ we fix $AB=10$ and $BC=4.$ Without the loss of generality, we consider $C$ on only one side of $\overline{AB}.$

 As shown below, all locations for $C$ at which $\triangle ABC$ is an obtuse triangle are indicated in red, excluding the endpoints.
                    \begin{figure}[H]
                        \centering
                        \includegraphics[width=0.5\textwidth]{Solution/AIME_2021_II_5_Solution_Figure_2.png}
                    \end{figure}
                    Note that:

 <ol style="margin-left: 1.5em;"> \item The region in which $\angle B$ is obtuse is determined by construction. 

 \item The region in which $\angle C$ is obtuse is determined by the corollaries of the Inscribed Angle Theorem.

 </ol>For any fixed value of $s,$ the height from $C$ is fixed. We need obtuse $\triangle ABC$ to be unique, so there can only be one possible location for $C.$ As shown below, all possible locations for $C$ are on minor arc $\widehat{C_1C_2},$ including $C_1$ but excluding $C_2.$
                    \begin{figure}[H]
                        \centering
                        \includegraphics[width=0.5\textwidth]{Solution/AIME_2021_II_5_Solution_Figure_3.png}
                    \end{figure}
                    Let the brackets denote areas: 

 <ul style="list-style-type:square;"> \item If $C=C_1,$ then $[ABC]$ will be minimized (attainable). By the same base and height and the Inscribed Angle Theorem, we have
 \begin{align*} [ABC]&=[ABD] \\ &=\frac12\cdot BD\cdot DA \\ &=\frac12\cdot BD\cdot \sqrt{AB^2-BD^2} \\ &=\frac12\cdot 4\cdot \sqrt{10^2-4^2} \\ &=2\sqrt{84}. \end{align*}
 

 \item If $C=C_2,$ then $[ABC]$ will be maximized (unattainable). For this right triangle, we have
 \begin{align*} [ABC]&=\frac12\cdot AB\cdot BC \\ &=\frac12\cdot 10\cdot 4 \\ &=20. \end{align*}


 </ul>Finally, the set of all such $s$ is $[a,b)=\left[2\sqrt{84},20\right),$ from which $a^2+b^2=\boxed{736}.$

 





\textbf{Solution 6}

Let a triangle in $\tau(s)$ be $ABC$, where $AB = 4$ and $BC = 10$. We will proceed with two cases:

 Case 1: $\angle ABC$ is obtuse. If $\angle ABC$ is obtuse, then, if we imagine $AB$ as the base of our triangle, the height can be anything in the range $(0,10)$; therefore, the area of the triangle will fall in the range of $(0, 20)$.

 Case 2: $\angle BAC$ is obtuse. Then, if we imagine $AB$ as the base of our triangle, the height can be anything in the range $\left(0, \sqrt{10^{2} - 4^{2}}\right)$. Therefore, the area of the triangle will fall in the range of $\left(0, 2 \sqrt{84}\right)$.

 If $s < 2 \sqrt{84}$, there will exist two types of triangles in $\tau(s)$ - one type with $\angle ABC$ obtuse; the other type with $\angle BAC$ obtuse. If $s \geq 2 \sqrt{84}$, as we just found, $\angle BAC$ cannot be obtuse, so therefore, there is only one type of triangle - the one in which $\angle ABC$ is obtuse. Also, if $s > 20$, no triangle exists with lengths $4$ and $10$. Therefore, $s$ is in the range $\left[ 2 \sqrt{84}, 20\right)$, so our answer is $\left(2 \sqrt{84}\right)^{2} + 20^{2} = \boxed{736}$.

 Alternatively, refer to Solution 5 for the geometric interpretation.

 





\textbf{Solution 7}


                    \begin{figure}[H]
                        \centering
                        \includegraphics[width=0.5\textwidth]{Solution/AIME_2021_II_5_Solution_Figure_4.png}
                    \end{figure}
                    Let's rephrase the condition. It is required to find such values of the area of an obtuse triangle with sides $4$ and $10,$ when there is exactly one such obtuse triangle. In the diagram, $AB = 4, AC = 10.$

 The largest area of triangle with sides $4$ and $10$ is $20$ for a right triangle with legs $4$ and $10$ ($AC\perp AB$).

 The diagram shows triangles with equal heights. The yellow triangle $ABC'$ has the longest side $BC',$ the blue triangle $ABC$ has the longest side $AC.$If $BC\perp AB,$ then$BC = \sqrt {AC^2 – AB^2} = 2 \sqrt{21}$ the area is equal to $4\sqrt{21}.$ In the interval, the blue triangle $ABC$ is acute-angled, the yellow triangle $ABC'$ is obtuse-angled. Their heights and areas are equal. The condition is met.

 If the area is less than $4\sqrt{21},$ both triangles are obtuse, not equal, so the condition is not met.

 Therefore, $s$ is in the range $\left[ 4 \sqrt{21}, 20\right)$, so answer is $\left(4 \sqrt{21}\right)^{2} + 20^{2} = \boxed{736}$

 \textbf{vladimir.shelomovskii@gmail.com, vvsss}

 





\textbf{Solution 8}

If $4$ and $10$ are the shortest sides and $\angle C$ is the included angle, then the area is 
 \[\frac{4\cdot10\cdot\sin\angle C}{2} = 20\sin\angle C.\]Because $0\leq\sin\angle C\leq1$, the maximum value of $20\sin\angle C$ is $20$, so $s\leq20$.

 If $4$ is a shortest side and $10$ is the longest side, the length of the other short side is $4\cos\angle C+2\sqrt{4\cos^2 \angle C+21}$ by law of cosines, and the area is $2\left(4\cos\angle C+2\sqrt{4\cos^2\angle C+21}\right)\sqrt{1-\cos\angle C}$. Because $-1\le \cos\angle C\le 0$, this is minimized if $\cos\angle C=0$, where $s=4\sqrt{21}$.

 So, the answer is $20^2+\left(4\sqrt{21}\right)^2=\boxed{736}$.

 




	\item (\textit{AIME II 2021}) For any finite set $S$, let $|S|$ denote the number of elements in $S$. Find the number of ordered pairs $(A,B)$ such that $A$ and $B$ are (not necessarily distinct) subsets of $\{1,2,3,4,5\}$ that satisfy

 \[|A| \cdot |B| = |A \cap B| \cdot |A \cup B|\]


 
 \textbf{Answer: }454



\textbf{Solution 1}

By PIE, $|A|+|B|-|A \cap B| = |A \cup B|$. Substituting into the equation and factoring, we get that $(|A| - |A \cap B|)(|B| - |A \cap B|) = 0$, so therefore $A \subseteq B$ or $B \subseteq A$. WLOG $A\subseteq B$, then for each element there are $3$ possibilities, either it is in both $A$ and $B$, it is in $B$ but not $A$, or it is in neither $A$ nor $B$. This gives us $3^{5}$ possibilities, and we multiply by $2$ since it could have also been the other way around. Now we need to subtract the overlaps where $A=B$, and this case has $2^{5}=32$ ways that could happen. It is $32$ because each number could be in the subset or it could not be in the subset. So the final answer is $2\cdot 3^5 - 2^5 = \boxed{454}$.

 





\textbf{Solution 2}

We denote $\Omega = \left\{ 1 , 2 , 3 , 4 , 5 \right\}$.We denote $X = A \cap B$, $Y = A \backslash \left( A \cap B \right)$, $Z = B \backslash \left( A \cap B \right)$, $W = \Omega \backslash \left( A \cup B \right)$.

 Therefore, $X \cup Y \cup Z \cup W = \Omega$ and the intersection of any two out of sets $X$, $Y$, $Z$, $W$ is an empty set.Therefore, $\left( X , Y , Z , W \right)$ is a partition of $\Omega$.

 Following from our definition of $X$, $Y$, $Z$, we have $A \cup B = X \cup Y \cup Z$.

 Therefore, the equation 

 
 \[|A| \cdot |B| = |A \cap B| \cdot |A \cup B|\]

 can be equivalently written as

 
 \[\left( | X | + | Y | \right) \left( | X | + | Z | \right) = | X | \left( | X | + | Y | + | Z | \right) .\]

 This equality can be simplified as

 
 \[| Y | \cdot | Z | = 0 .\]

 Therefore, we have the following three cases: (1) $|Y| = 0$ and $|Z| \neq 0$, (2) $|Z| = 0$ and $|Y| \neq 0$, (3) $|Y| = |Z| = 0$.Next, we analyze each of these cases, separately.

 Case 1: $|Y| = 0$ and $|Z| \neq 0$.

 In this case, to count the number of solutions, we do the complementary counting. 

 First, we count the number of solutions that satisfy $|Y| = 0$.

 Hence, each number in $\Omega$ falls into exactly one out of these three sets: $X$, $Z$, $W$.Following from the rule of product, the number of solutions is $3^5$.

 Second, we count the number of solutions that satisfy $|Y| = 0$ and $|Z| = 0$.

 Hence, each number in $\Omega$ falls into exactly one out of these two sets: $X$, $W$.Following from the rule of product, the number of solutions is $2^5$.

 Therefore, following from the complementary counting, the number of solutions in this case is equal to the number of solutions that satisfy $|Y| = 0$ minus the number of solutions that satisfy $|Y| = 0$ and $|Z| = 0$, i.e., $3^5 - 2^5$.

 Case 2: $|Z| = 0$ and $|Y| \neq 0$. 

 This case is symmetric to Case 1. Therefore, the number of solutions in this case is the same as the number of solutions in Case 1, i.e., $3^5 - 2^5$.

 Case 3: $|Y| = 0$ and $|Z| = 0$.

 Recall that this is one part of our analysis in Case 1. Hence, the number solutions in this case is $2^5$.

 By putting all cases together, following from the rule of sum, the total number of solutions is equal to

 
 \begin{align*} \left( 3^5 - 2^5 \right) + \left( 3^5 - 2^5 \right) + 2^5 & = 2 \cdot 3^5 - 2^5 \\ & = \boxed{454} . \end{align*}


 





\textbf{Solution 3 (Principle of Inclusion-Exclusion)}

By the \textbf{Principle of Inclusion-Exclusion (abbreviated as PIE)}, we have $|A \cup B|=|A|+|B|-|A \cap B|,$ from which we rewrite the given equation as 
 \[|A| \cdot |B| = |A \cap B| \cdot \left(|A|+|B|-|A \cap B|\right).\]Rearranging and applying Simon's Favorite Factoring Trick give
 \begin{align*} |A| \cdot |B| &= |A \cap B|\cdot|A| + |A \cap B|\cdot|B| - |A \cap B|^2 \\ |A| \cdot |B| - |A \cap B|\cdot|A| - |A \cap B|\cdot|B| &= - |A \cap B|^2 \\ \left(|A| - |A \cap B|\right)\cdot\left(|B| - |A \cap B|\right) &=0, \end{align*}
from which at least one of the following is true: 

 
                \begin{itemize}
                    \item $|A|=|A \cap B|$
                \end{itemize}


                
                \begin{itemize}
                    \item $|B|=|A \cap B|$
                \end{itemize}


                Let $|A \cap B|=k.$ For each value of $k\in\{0,1,2,3,4,5\},$ we will use PIE to count the ordered pairs $(A,B):$

 Suppose $|A|=k.$ There are $\binom{5}{k}$ ways to choose the elements for $A.$ These $k$ elements must also appear in $B.$ Next, there are $2^{5-k}$ ways to add any number of the remaining $5-k$ elements to $B$ (Each element has $2$ options: in $B$ or not in $B.$). There are $\binom{5}{k}2^{5-k}$ ordered pairs for $|A|=k.$ Similarly, there are $\binom{5}{k}2^{5-k}$ ordered pairs for $|B|=k.$ 

 To fix the overcount, we subtract the number of ordered pairs that are counted twice, in which $|A|=|B|=k.$ There are $\binom{5}{k}$ such ordered pairs.

 Therefore, there are 
 \[2\binom{5}{k}2^{5-k}-\binom{5}{k}\] ordered pairs for $|A \cap B|=k.$

 Two solutions follow from here:

 





\textbf{Solution 3.1 (Binomial Theorem)}

The answer is 
 \begin{align*} \sum_{k=0}^{5}\left[2\binom{5}{k}2^{5-k}-\binom{5}{k}\right] &= 2\sum_{k=0}^{5}\binom{5}{k}2^{5-k}-\sum_{k=0}^{5}\binom{5}{k} \\ &=2(2+1)^5-(1+1)^5 \\ &=2(243)-32 \\ &=\boxed{454}. \end{align*}
~MRENTHUSIASM

 





\textbf{Solution 3.2 (Bash)}

The answer is 
 \begin{align*} &\hspace{5.125mm}\sum_{k=0}^{5}\left[2\binom{5}{k}2^{5-k}-\binom{5}{k}\right] \\ &=\left[2\binom{5}{0}2^{5-0}-\binom{5}{0}\right] + \left[2\binom{5}{1}2^{5-1}-\binom{5}{1}\right] + \left[2\binom{5}{2}2^{5-2}-\binom{5}{2}\right] + \left[2\binom{5}{3}2^{5-3}-\binom{5}{3}\right] + \left[2\binom{5}{4}2^{5-4}-\binom{5}{4}\right] + \left[2\binom{5}{5}2^{5-5}-\binom{5}{5}\right] \\ &=\left[2\left(1\right)2^5-1\right] + \left[2\left(5\right)2^4-5\right] + \left[2\left(10\right)2^3-10\right] + \left[2\left(10\right)2^2-10\right] + \left[2\left(5\right)2^1-5\right] + \left[2\left(1\right)2^0-1\right] \\ &=63+155+150+70+15+1 \\ &=\boxed{454}. \end{align*}
~MRENTHUSIASM

 





\textbf{Solution 4 (Simple Bash)}

Proceed with Solution 1 to get $(|A| - |A \cap B|)(|B| - |A \cap B|) = 0$. WLOG, assume $|A| = |A \cap B|$. Thus, $A \subseteq B$.

 Since $A \subseteq B$, if $|B| = n$, there are $2^n$ possible sets $A$, and there are also ${5 \choose n}$ ways of choosing such $B$.

 Therefore, the number of possible pairs of sets $(A, B)$ is

 
 \[\sum_{k=0}^{5} 2^n {5 \choose n}\]

 We can compute this manually since it's only from $k=0$ to $5$, and computing gives us $243$. We can double this result for $B \subseteq A$, and we get $2(243) = 486$.

 However, we have double counted the cases where $A$ and $B$ are the same sets. There are $32$ possible such cases, so we subtract $32$ from $486$ to get $\boxed{454}$.

 




	\item (\textit{AIME II 2021}) Let $a, b, c,$ and $d$ be real numbers that satisfy the system of equations

 \begin{align*} a + b &= -3, \\ ab + bc + ca &= -4, \\ abc + bcd + cda + dab &= 14, \\ abcd &= 30. \end{align*}

There exist relatively prime positive integers $m$ and $n$ such that

 \[a^2 + b^2 + c^2 + d^2 = \frac{m}{n}.\]Find $m + n$.


 
 \textbf{Answer: }145



\textbf{Solution 1}

From the fourth equation we get $d=\frac{30}{abc}.$ Substitute this into the third equation and you get $abc + \frac{30(ab + bc + ca)}{abc} = abc - \frac{120}{abc} = 14$. Hence $(abc)^2 - 14(abc)-120 = 0$. Solving, we get $abc = -6$ or $abc = 20$. From the first and second equation, we get $ab + bc + ca = ab-3c = -4 \Longrightarrow ab = 3c-4$. If $abc=-6$, substituting we get $c(3c-4)=-6$. If you try solving this you see that this does not have real solutions in $c$, so $abc$ must be $20$. So $d=\frac{3}{2}$. Since $c(3c-4)=20$, $c=-2$ or $c=\frac{10}{3}$. If $c=\frac{10}{3}$, then the system $a+b=-3$ and $ab = 6$ does not give you real solutions. So $c=-2$. Since you already know $d=\frac{3}{2}$ and $c=-2$, so you can solve for $a$ and $b$ pretty easily and see that $a^{2}+b^{2}+c^{2}+d^{2}=\frac{141}{4}$. So the answer is $\boxed{145}$.

 





\textbf{Solution 2}

Note that $ab + bc + ca = -4$ can be rewritten as $ab + c(a+b) = -4$. Hence, $ab = 3c - 4$.

 Rewriting $abc+bcd+cda+dab = 14$, we get $ab(c+d) + cd(a+b) = 14$. Substitute $ab = 3c - 4$ and solving, we get 
 \[3c^{2} - 4c - 4d - 14 = 0.\] We refer to this as Equation 1.

 Note that $abcd = 30$ gives $(3c-4)cd = 30$. So, $3c^{2}d - 4cd = 30$, which implies $d(3c^{2} - 4c) = 30$ or 
 \[3c^{2} - 4c = \frac{30}{d}.\] We refer to this as Equation 2.

 Substituting Equation 2 into Equation 1 gives, $\frac{30}{d} - 4d - 14 = 0$.

 Solving this quadratic yields that $d \in \left\{-5, \frac{3}{2}\right\}$.

 Now we just try these two cases:

 For $d = \frac{3}{2}$ substituting in Equation 1 gives a quadratic in $c$ which has roots $c \in \left\{\frac{10}{3}, -2\right\}$.

 Again trying cases, by letting $c = -2$, we get $ab = 3c-4$, Hence $ab = -10$.We know that $a + b = -3$, Solving these we get $a = -5, b = 2$ or $a= 2, b = -5$ (doesn't matter due to symmetry in $a,b$).

 So, this case yields solutions $(a,b,c,d) = \left(-5, 2 , -2, \frac{3}{2}\right)$.

 Similarly trying other three cases, we get no more solutions, Hence this is the solution for $(a,b,c,d)$.

 Finally, $a^{2} + b^{2} + c^{2} + d^{2} = 25 + 4 + 4 + \frac{9}{4} = \frac{141}{4} = \frac{m}{n}$.

 Therefore, $m + n = 141 + 4 = \boxed{145}$.

 





\textbf{Solution 3}

For simplicity purposes, we number the given equations $(1),(2),(3),$ and $(4),$ in that order.

 Rearranging $(2)$ and solving for $c,$ we have
 \begin{align*} ab+(a+b)c&=-4 \\ ab-3c&=-4 \\ c&=\frac{ab+4}{3}. \hspace{14mm} (5) \end{align*}
Substituting $(5)$ into $(4)$ and solving for $d,$ we get
 \begin{align*} ab\left(\frac{ab+4}{3}\right)d&=30 \\ d&=\frac{90}{ab(ab+4)}. \hspace{5mm} (6)  \end{align*}
Substituting $(5)$ and $(6)$ into $(3)$ and simplifying, we rewrite the left side of $(3)$ in terms of $a$ and $b$ only:
 \begin{align*} ab\left[\frac{ab+4}{3}\right] + b\left[\frac{ab+4}{3}\right]\left[\frac{90}{ab(ab+4)}\right] + \left[\frac{ab+4}{3}\right]\left[\frac{90}{ab(ab+4)}\right]a + \left[\frac{90}{ab(ab+4)}\right]ab &= 14 \\ ab\left[\frac{ab+4}{3}\right] + \underbrace{\frac{30}{a} + \frac{30}{b}}_{\text{Group them.}} + \frac{90}{ab+4} &= 14 \\ ab\left[\frac{ab+4}{3}\right] + \frac{30(a+b)}{ab} + \frac{90}{ab+4} &= 14 \\ ab\left[\frac{ab+4}{3}\right] + \underbrace{\frac{-90}{ab} + \frac{90}{ab+4}}_{\text{Group them.}} &= 14 \\ ab\left[\frac{ab+4}{3}\right] - \frac{360}{ab(ab+4)}&=14. \end{align*}
Let $t=ab(ab+4),$ from which 
 \[\frac{t}{3}-\frac{360}{t}=14.\] Multiplying both sides by $3t,$ rearranging, and factoring give $(t+18)(t-60)=0.$ Substituting back and completing the squares produce
 \begin{align*} \left[ab(ab+4)+18\right]\left[ab(ab+4)-60\right]&=0 \\ \left[(ab)^2+4ab+18\right]\left[(ab)^2+4ab-60\right]&=0 \\ \underbrace{\left[(ab+2)^2+14\right]}_{ab+2=\pm\sqrt{14}i\implies ab\not\in\mathbb R}\underbrace{\left[(ab+2)^2-64\right]}_{ab+2=\pm8}&=0 \\ ab&=6,-10. \end{align*}
If $ab=6,$ then combining this with $(1),$ we know that $a$ and $b$ are the solutions of the quadratic $x^2+3x+6=0.$ Since the discriminant is negative, neither $a$ nor $b$ is a real number.

 If $ab=-10,$ then combining this with $(1),$ we know that $a$ and $b$ are the solutions of the quadratic $x^2+3x-10=0,$ or $(x+5)(x-2)=0,$ from which $\{a,b\}=\{-5,2\}.$ Substituting $ab=-10$ into $(5)$ and $(6),$ we obtain $c=-2$ and $d=\frac32,$ respectively. Together, we have 
 \[a^2+b^2+c^2+d^2=\frac{141}{4},\] so the answer is $141+4=\boxed{145}.$

 





\textbf{Solution 4 (Way Too Long)}

Let the four equations from top to bottom be listed 1 through 4 respectively. We factor equation 3 like so: 
 \[abc+d(ab+bc+ca)=14\]Then we plug in equation 2 to receive $abc-4d=14$. By equation 4 we get $abc=\frac{30}{d}$. Plugging in, we get $\frac{30}{d}-4d=14$. Multiply by $d$ on both sides to get the quadratic equation $4d^2+14d-30=0$. Solving using the quadratic equation, we receive $d=\frac{3}{2},d=-5$. So, we have to test which one is correct. We repeat a similar process as we did above for equations 1 and 2. We factor equation 2 to get 
 \[ab+c(a+b)=-4\]After plugging in equation 1, we get $ab-3c=-4$. Now we convert it into a quadratic to receive $3c^2-4c-abc=0$. The value of $abc$ will depend on $d$. So we obtain the discriminant $16+12abc$.Let $d = -5$.Then $abc = \frac{30}{-5}$, so $abc=-6$, discriminant is $16-72$, which makes this a dead end. Thus $d=\frac{3}{2}$For $d=\frac{3}{2}$, making $abc=20$. This means the discriminant is just $256$, so we obtain two values for $c$ as well. We get either $c=\frac{10}{3}$ or $c=-2$. So, we must AGAIN test which one is correct.We know $ab=3c-4$, and $a+b=-3$, so we use these values for testing.Let $c=\frac{10}{3}$.Then $ab=6$, so $a=\frac{6}{b}$. We thus get $\frac{6}{b}+b=-3$, which leads to the quadratic $b^2+3b+6$. The discriminant for this is $9-24$. That means this value of $c$ is wrong, so $c=-2$. Thus we get polynomial $b^2+3b-10$. The discriminant this time is $49$, so we get two values for $b$. Through simple inspection, you may see they are interchangeable, as if you take the value $b=2$, you get $a=-5$. If you take the value $b=-5$, you get $a=2$. So it doesn't matter. That means the sum of all their squares is
 \[\frac{9}{4}+4+4+25=\frac{141}{4},\]so the answer is $141+4=\boxed{145}.$

 





\textbf{Solution 5}

Let the four equations from top to bottom be listed $(1)$ through $(4)$ respectively. Multiplying both sides of $(3)$ by $d$ and factoring some terms gives us $abcd + d^2(ab + ac + bc) = 14d$. Substituting using equations $(4)$ and $(2)$ gives us $30 -4 d^2 = 14d$, and solving gives us $d = -5$ or $d = \frac{3}{2}$. Plugging this back into $(3)$ gives us $abc + d(ab + ac + bc) = abc + (-5)(-4) = abc + 20 = 14$, or using the other solution for $d$ gives us $abc - 6 = 14$. Solving both of these equations gives us $abc = -6$ when $d = -5$ and $abc = 20$ when $d = \frac{3}{2}$. 

 Multiplying both sides of $(2)$ by $c$ and factoring some terms gives us $abc + c^2 (a + b) = abc -3c^2 = -4c$. Testing $abc = -6$ will give us an imaginary solution for $c$, so therefore $abc = 20$ and $d = \frac{3}{2}$. This gets us $20 - 3c^2 = -4c$. Solving for $c$ gives us $c = \frac{3}{10}$ or $c = -2$. With a bit of testing, we can see that the correct value of $c$ is $c=-2$. Now we know $a+b = -3$ and $ab + bc + ca = ab + c(a+b) = ab + 6 = -4$, $ab = -10$, and it is obvious that $a = -5$ and $b = 2$ or the other way around, and therefore, $a^2 + b^2 + c^2 + d^2 = 25 + 4 + 4 + \frac{9}{4} = \frac{141}{4}$, giving us the answer $141 + 4 = \boxed{145}$.

 




	\item (\textit{AIME II 2021}) An ant makes a sequence of moves on a cube where a move consists of walking from one vertex to an adjacent vertex along an edge of the cube. Initially the ant is at a vertex of the bottom face of the cube and chooses one of the three adjacent vertices to move to as its first move. For all moves after the first move, the ant does not return to its previous vertex, but chooses to move to one of the other two adjacent vertices. All choices are selected at random so that each of the possible moves is equally likely. The probability that after exactly $8$ moves that ant is at a vertex of the top face on the cube is $\frac{m}{n}$, where $m$ and $n$ are relatively prime positive integers. Find $m + n.$


 
 \textbf{Answer: }049



\textbf{Solution 1 (Four-Variable Recursion)}

For all positive integers $k,$ let

 
                \begin{itemize}
                    \item $N(k,\mathrm{BB})$ be the number of ways to make a sequence of exactly $k$ moves, where the last move is from the bottom face to the bottom face.
                \end{itemize}


                
                \begin{itemize}
                    \item $N(k,\mathrm{BT})$ be the number of ways to make a sequence of exactly $k$ moves, where the last move is from the bottom face to the top face.
                \end{itemize}


                
                \begin{itemize}
                    \item $N(k,\mathrm{TB})$ be the number of ways to make a sequence of exactly $k$ moves, where the last move is from the top face to the bottom face.
                \end{itemize}


                
                \begin{itemize}
                    \item $N(k,\mathrm{TT})$ be the number of ways to make a sequence of exactly $k$ moves, where the last move is from the top face to the top face.
                \end{itemize}


                The base case occurs at $k=1,$ from which $\left(N(1,\mathrm{BB}),N(1,\mathrm{BT}),N(1,\mathrm{TB}),N(1,\mathrm{TT})\right)=(2,1,0,0).$ 

 Suppose the ant makes exactly $k$ moves for some $k\geq2.$ We perform casework on its last move:

 <ol style="margin-left: 1.5em;"> \item If its last move is from the bottom face to the bottom face, then its next move has

 <ul style="margin-left: 1.5em; list-style-type: circle"> \item $1$ way to move from the bottom face to the bottom face. 

 \item $1$ way to move from the bottom face to the top face.

 </ul> \item If its last move is from the bottom face to the top face, then its next move has $2$ ways to move from the top face to the top face. 

 \item If its last move is from the top face to the bottom face, then its next move has $2$ ways to move from the bottom face to the bottom face. 

 <li>If its last move is from the top face to the top face, then its next move has

 <ul style="margin-left: 1.5em; list-style-type: circle"> \item $1$ way to move from the top face to the bottom face. 

 \item $1$ way to move from the top face to the top face.

 </ul></li></ol>Alternatively, this recursion argument is illustrated below, where each dashed arrow indicates $1$ way, and each solid arrow indicates $2$ ways:
                    \begin{figure}[H]
                        \centering
                        \includegraphics[width=0.5\textwidth]{Solution/AIME_2021_II_8_Solution_Figure_1.png}
                    \end{figure}
                    Therefore, we have the following relationships:
 \begin{align*} N(1,\mathrm{BB})&=2, \\  N(1,\mathrm{BT})&=1, \\ N(1,\mathrm{TB})&=0, \\ N(1,\mathrm{TT})&=0, \\ N(k+1,\mathrm{BB})&=N(k,\mathrm{BB})+2\cdot N(k,\mathrm{TB}), \\ N(k+1,\mathrm{BT})&=N(k,\mathrm{BB}), \\ N(k+1,\mathrm{TB})&=N(k,\mathrm{TT}), \\ N(k+1,\mathrm{TT})&=N(k,\mathrm{TT})+2\cdot N(k,\mathrm{BT}). \end{align*}
Using these equations, we recursively fill out the table below:
 \[\begin{array}{c||c|c|c|c|c|c|c|c}    \hspace{7mm}&\hspace{6.5mm}&\hspace{6.5mm}&\hspace{6.75mm}&\hspace{6.75mm}&\hspace{6.75mm}&\hspace{6.75mm}&& \\ [-2.5ex] \boldsymbol{k} & \boldsymbol{1} & \boldsymbol{2} & \boldsymbol{3} & \boldsymbol{4} & \boldsymbol{5} & \boldsymbol{6} & \boldsymbol{7} & \boldsymbol{8} \\  \hline \hline &&&&&&&& \\ [-2.25ex] \boldsymbol{N(k,\mathrm{BB})} &2&2&2&6&18&38&66&118 \\ \hline   &&&&&&&& \\ [-2.25ex] \boldsymbol{N(k,\mathrm{BT})} &1&2&2&2&6&18&38&66 \\ \hline  &&&&&&&& \\ [-2.25ex] \boldsymbol{N(k,\mathrm{TB})} &0&0&2&6&10&14&26&62 \\ \hline    &&&&&&&& \\ [-2.25ex] \boldsymbol{N(k,\mathrm{TT})} &0&2&6&10&14&26&62&138 \\ \hline \hline &&&&&&&& \\ [-2.25ex] \textbf{Total}&\boldsymbol{3}&\boldsymbol{6}&\boldsymbol{12}&\boldsymbol{24}&\boldsymbol{48}&\boldsymbol{96}&\boldsymbol{192}&\boldsymbol{384} \end{array}\]By the Multiplication Principle, there are $3\cdot2^{k-1}$ ways to make exactly $k$ moves. So, we must get 
 \[N(k,\mathrm{BB})+N(k,\mathrm{BT})+N(k,\mathrm{TB})+N(k,\mathrm{TT})=3\cdot2^{k-1}\] for all values of $k.$

 Finally, the requested probability is 
 \[\frac{N(8,\mathrm{BT})+N(8,\mathrm{TT})}{N(8,\mathrm{BB})+N(8,\mathrm{BT})+N(8,\mathrm{TB})+N(8,\mathrm{TT})}=\frac{66+138}{118+66+62+138}=\frac{204}{384}=\frac{17}{32},\] from which the answer is $17+32=\boxed{049}.$

 





\textbf{Solution 2 (Markov Chain and Dynamic Programming)}

Let the state from bottom to top be $B2T,$ from top to top be $T2T,$ from top to bottom be $T2B,$ and from bottom to bottom be $B2B.$ We can draw the following State Transition Diagram with \href{https://en.wikipedia.org/wiki/Markov\_chain}{Markov Chain}. The numbers on the transition arc are the transition probabilities. 

 
                    \begin{figure}[H]
                        \centering
                        \includegraphics[width=0.5\textwidth]{Solution/AIME_2021_II_8_Solution_Figure_2.png}
                    \end{figure}
                    The probabilities of being in a state after $n$ steps and after $n-1$ steps has the following relationships:
 \begin{align*} B2T(n) &= B2B(n-1) \cdot \frac12\\ T2T(n) &= B2T(n-1) + T2T(n-1) \cdot \frac12\\ T2B(n) &= T2T(n-1) \cdot \frac12\\ B2B(n) &= T2B(n-1) + B2B(n-1) \cdot \frac12 \end{align*}
Those probabilities are calculated by \href{https://en.wikipedia.org/wiki/Dynamic\_programming}{Dynamic Programming} in the following table:
 \[\begin{array}{c|cccc}  &  &  &  & \\ [-2ex] n & B2T(n) & T2T(n) & T2B(n) & B2B(n) \\ [1ex] \hline  &  &  &  & \\ [-1ex] 1 & \frac13 & 0 & 0 & \frac23\\  &  &  &  & \\ 2 & \frac23 \cdot \frac12 = \frac13 & \frac13 & 0 & \frac23 \cdot \frac12 = \frac13 \\  &  &  &  & \\ 3 & \frac13 \cdot \frac12 = \frac16 & \frac13 + \frac13 \cdot \frac12 = \frac12 & \frac13 \cdot \frac12 = \frac16 & \frac13 \cdot \frac12 = \frac16 \\  &  &  &  & \\ 4 & \frac16 \cdot \frac12 = \frac{1}{12} & \frac16 + \frac12 \cdot \frac12 = \frac{5}{12} & \frac12 \cdot \frac12 = \frac14 & \frac16 + \frac16 \cdot \frac12 = \frac14 \\  &  &  &  & \\ 5 & \frac14 \cdot \frac12 = \frac18 & \frac{1}{12} + \frac{5}{12} \cdot \frac{1}{2} = \frac{7}{24} & \frac{5}{12} \cdot \frac12 = \frac{5}{24} & \frac14 + \frac14 \cdot \frac12 = \frac38 \\  &  &  &  & \\ 6 & \frac38 \cdot \frac12 = \frac{3}{16} & \frac18 + \frac{7}{24} \cdot \frac12 = \frac{13}{48} & \frac{7}{24} \cdot \frac12 = \frac{7}{48} & \frac{5}{24} + \frac38 \cdot \frac12 = \frac{19}{48} \\  &  &  &  & \\ 7 & \frac{19}{48} \cdot \frac12 = \frac{19}{96}&  \frac{3}{16} + \frac{13}{48} \cdot \frac{1}{2} = \frac{31}{96}& \frac{13}{48} \cdot \frac{1}{2} = \frac{13}{96} & \frac{7}{48} + \frac{19}{48} \cdot \frac12 = \frac{11}{32}\\  &  &  &  & \\ 8 & \frac{11}{32} \cdot \frac{1}{2} = \frac{11}{64} & \frac{19}{96} + \frac{31}{96} \cdot \frac12 = \frac{23}{64} & \frac{31}{96} \cdot \frac12 = \frac{31}{192} & \frac{13}{96} + \frac{11}{32} \cdot \frac12 = \frac{59}{192} \\ [1ex] \end{array}\]Finally, the requested probability is $\frac{11}{64} + \frac{23}{64} = \frac{17}{32},$ from which the answer is $17 + 32 = \boxed{049}.$

 





\textbf{Solution 3 (One-Variable Recursion)}

Note that we don't care which exact vertex the ant is located at, just which level (either top face or bottom face).Consider the ant to be on any of the two levels and having moved at least one move. Define $p_i$ to be the probability that after $i$ moves, the ant ends up on the level it started on.

 On the first move, the ant can stay on the bottom level with $\frac{2}{3}$ chance and $7$ moves left. Or, it can move to the top level with $\frac{1}{3}$ chance and $6$ moves left (it has to spend another on the top as it can not return immediately). So the requested probability is $P = \frac{2}{3}(1 - p_7) + \frac{1}{3}p_6$. 

 Consider when the ant has $i$ moves left (and it's not the ant's first move). It can either stay on its current level with $\frac{1}{2}$ chance and $i - 1$ moves left, or travel to the opposite level with $\frac{1}{2}$ chance, then move to another vertex on the opposite level, to have $i - 2$ moves left. Thus we obtain the recurrence
 \[p_i = \frac{1}{2}p_{i - 1} + \frac{1}{2}(1 - p_{i - 2})\]

 Computing $p_i$ with the starting conditions $p_0 = 1$ and $p_1 = \frac{1}{2}$, we obtain $p_6 = \frac{33}{64}$ and $p_7 = \frac{59}{128}$. Hence $P = \frac{2}{3}(1 - p_7) + \frac{1}{3}p_6= \frac{17}{32}$ as desired; $\boxed{049}$.

 





\textbf{Solution 4 (Casework)}

On each move, we can either stay on the level we previously were (stay on the bottom/top) or switch levels (go from top to bottom and vise versa). Since we start on the bottom, ending on the top means that we will have to switch an odd number of times; since we cannot switch twice in a row, over an eight-move period we can either make one or three switches. Furthermore, once we switch to a level we can choose one of two directions of traveling on that level: clockwise or counterclockwise (since we can't go back to our previous move, our first move on the level after switching determines our direction). 

 
            \begin{enumerate}
                <li>\textbf{Case 1: one switch.} Our one switch can either happen at the start/end of our moves, or in the middle. There are $4 + 24 = 28$ ways to do this, outlined below.
            \begin{enumerate}
                \item \textbf{Subcase 1: switch happens at ends.} If our first move is a switch, then there are two ways to determine the direction we travel along the top layer. Multiply by $2$ to count for symmetry (last move is a switch) so this case yields $2^2 = 4$ possibilities.\item \textbf{Subcase 2: switch happens in the middle.} There are six places for the switch to happen; the switch breaks the sequences of moves into two chains, with each having $2$ ways to choose their direction of travel. This case yields $6 \cdot 2^2 = 24$ possibilities.
            \end{enumerate}
            </li><li>\textbf{Case 2: three switches.} Either two, one, or none of our switches occur at the start/end of our moves. There are $16 + 96 + 64 = 176$ ways to do this, outlined below. (Keep in mind we can't have two switches in a row.)
            \begin{enumerate}
                \item \textbf{Subcase 1: start and end with a switch.} Since our third switch can't be in moves $2$ or $7$, there are four ways to place our switch, breaking our sequence into two chains. This case yields $4 \cdot 2^2 = 16$ possibilities.\item \textbf{Subcase 2: one of our switches is at the start/end.} WLOG our first move is a switch; moves $2$ and $8$ cannot be switches. We can choose $2$ from any of the remaining $5$ moves to be switches, but we have to subtract the $4$ illegal cases where the two switches are in a row (3-4, 4-5, 5-6, 6-7). These three switches break our sequence into three chains; accounting for symmetry this case yields $2\left(\binom{5}{2} - 4\right) \cdot 2^3 = 96$ possibilities.\item \textbf{Subcase 3: all our switches are in the middle.} We choose $3$ from any of the $6$ middle moves to be our switches, but have to subtract the cases where at least two of them are in a row. If at least two switches are in a row, there are five places for the group of $2$ and four places for the third switch; however this overcounts the case where all three are in a row, which has $4$ possibilities. These three switches break our sequence into four chains, so this case yields $\left(\binom{6}{3} - 5 \cdot 4 + 4\right) \cdot 2^4 = 64$ possibilities.
            \end{enumerate}
            </li>
            \end{enumerate}
            Our probability is then $\frac{176 + 28}{3 \cdot 2^7} = \frac{17}{32}$, so the answer is $17+32=\boxed{049}$.

 





\textbf{Solution 5 (One-Variable Recursion and Casework)}

Let $p(n)$ be the probability that we are on the top when you get to the $n$th move, and $p(n-1)$ and $p(n-2)$ be the probability that you are on the top when you get to the $(n-1)$th move and $(n-2)$th move respectively.

 Now you can do some recursion, splitting up into cases:

 Case 1: You are not on the top for the $(n-1)$th move or the $(n-2)$th move. In this case, it is a $1-p(n-1)$ chance that you were not on the top for the $(n-1)$th move, and a $1-p(n-2)$ that you are not on the top for the $(n-2)$th move. This leads you to a $\frac{1}{2}\cdot(1-p(n-1))\cdot(1-p(n-2))$ chance for the first case. (The $\frac12$ is there because of the fact that you can go up and you can also stay on the bottom side, as you cannot return.)

 Case 2: You are on the top for the $(n-1)$th move, but not on the top for the $(n-2)$th move. This leads you to a $p(n-1)\cdot(1-p(n-2))$ probability (no extra components because you cannot return).

 Case 3: You are on top for both the $(n-1)$th move and the $(n-2)$th move. This leads to a probability of $\frac{1}{2}\cdot p(n-1)\cdot p(n-2)$ (adding the extra $\frac12$ because you can either stay on the top or go down.

 As the 4th case requires you to go down then up, but you cannot retrace, there is no $4$th case.

 These cases, added up, lead you to 
 \[p(n)=\frac{1}{2}\cdot(1-p(n-1))(1-p(n-2))+p(n-1)(1-p(n-2))+\frac{1}{2}\cdot p(n-1)p(n-2).\] This can be further simplified down by expanding and combining like terms to 
 \[p(n)=\frac{1+p(n-1)-p(n-2)}{2}.\]Then we must find $p(1)$ and $p(2)$. p(1) is trivially $\frac{1}{3}$. You can find $p(2)$ using basic probability techniques that is left as an exercise to the reader to get $\frac{2}{3}$. In the end, you plug in to get
 \[p(1)=\frac{1}{3}, P(2)=\frac{2}{3}, P(3)=\frac{2}{3}, P(4)=\frac{1}{2}, P(5)=\frac{5}{12}, P(6)=\frac{11}{24}, P(7)=\frac{25}{48}, P(8)=\boxed{\frac{17}{32}}.\]Therefore, the answer is $17+32=\boxed{049}$.

 





\textbf{Solution 6 (Casework)}

Illustration as

 \begin{verbatim}
T3\_\_\_\_\_\_\_\_\_T4|\          |\| \         | \|  \        |  \|  T2\_\_\_\_\_\_\_\_\_\_\_T1|   |       |   |B3\_ | \_\_\_\_\_B4   |\   |        \  | \  |         \ |  \ |          \|   \B2\_\_\_\_\_\_\_\_\_\_B1
\end{verbatim}WLOG we assume the ant starts from B1. After 8 moves, it will arrive on one of T2, T4, B1 and B3. And apparently T2/T4/B1 has the same possibility. In the other word, if we can find out the possibililty of arrival of $P(B1)$, $P(T2)=P(T4)=P(B3)=\frac{1}{3} (1-P(B1))$

 After 8 moves, there are $3\times 2^7$ cases. Now we consider the cases the ant can arrive B1.

 Since the ant moves along the edges, there are 6 different possible moves: +x, -x, +y, -y, +z, -z. For each direction x/y/z, the ant will always move even times and always move positively first and negatively then. So we do not need to differ +x and -x (so does y and z). And since the ant can not move backward immediately, no two x(or y or z) should be consecutive. Now we consider the possible sequence of x/y/z of 8 moves.

 Choice 1: only one direction, impossible since no two consecutive x/y/z

 Choice 2: only two direction. For example, if there are only x and y, there are only two possible sequence

 $xyxyxyxy$ $yxyxyxyx$

 There are total 6 cases for this scenario.

 Choice 3: all three directions, consider there are 4 x's and 2 y's and 2 z'sThere are 6 arrangements for 2 y's and 2 z's

 $yyzz$ $yzyz$ $yzzy$ $zyyz$ $zyzy$ $zzyy$

 Corresponding for the 4 x's placement numbers are 3, 5, 4, 4, 5, 3. Consider x/y/z to be rotational. The total is $(3+5+4)\times 2\times 3=72$.

 Therefore $P(B1)=\frac{(6+72)}{3*2^7}=\frac{13}{64}$

 Which gives us $P(T2)+P(T4)=\frac{2}{3} (1-P(B1))=\frac{17}{32}$

 So the answer is $17+32=\boxed{049}$.

 \textbf{Remark (Markov Chain)} This problem is similar to the following problems:

 
                \begin{itemize}
                    \item \href{https://artofproblemsolving.com/wiki/index.php/1985\_AIME\_Problems/Problem\_12}{1985 AIME Problem 12}
                \end{itemize}


                
                \begin{itemize}
                    \item \href{https://artofproblemsolving.com/wiki/index.php/2003\_AIME\_II\_Problems/Problem\_13}{2003 AIME II Problem 13}
                \end{itemize}


                
                \begin{itemize}
                    \item \href{https://artofproblemsolving.com/wiki/index.php/2021\_AMC\_12A\_Problems/Problem\_23}{2021 AMC 12A Problem 23}
                \end{itemize}


                
                \begin{itemize}
                    \item \href{https://artofproblemsolving.com/wiki/index.php/2022\_AMC\_8\_Problems/Problem\_25}{2022 AMC 8 Problem 25}
                \end{itemize}


                They can all be solved by \href{https://en.wikipedia.org/wiki/Markov\_chain}{Markov Chain} and \href{https://en.wikipedia.org/wiki/Dynamic\_programming}{Dynamic Programming}.

 Let $p_{ij} = P(X_{n+1} = j | X_n = i)$ be the probability that state $i$ transits to state $j$ on the next step, and $s_i$ be the probability of being in state $i$. It follows that $s_j$ is the sum of the products of $s_i$ and $p_{ij}$ of all the previous state $i$: 
 \[s_j = \sum_{i} (s_{i} \cdot p_{ij}).\]This formula can also be applied to solve \href{https://artofproblemsolving.com/wiki/index.php/2019\_AMC\_10B\_Problems/Problem\_22}{2019 AMC10B Problem 22}.

 




	\item (\textit{AIME II 2021}) Find the number of ordered pairs $(m, n)$ such that $m$ and $n$ are positive integers in the set $\{1, 2, ..., 30\}$ and the greatest common divisor of $2^m + 1$ and $2^n - 1$ is not $1$.


 
 \textbf{Answer: }295



\textbf{Solution 1}

This solution refers to the \textbf{Remarks} section.

 By the Euclidean Algorithm, we have 
 \[\gcd\left(2^m+1,2^m-1\right)=\gcd\left(2,2^m-1\right)=1.\]We are given that $\gcd\left(2^m+1,2^n-1\right)>1.$ Multiplying both sides by $\gcd\left(2^m-1,2^n-1\right)$ gives
 \begin{align*} \gcd\left(2^m+1,2^n-1\right)\cdot\gcd\left(2^m-1,2^n-1\right)&>\gcd\left(2^m-1,2^n-1\right) \\ \gcd\left(\left(2^m+1\right)\left(2^m-1\right),2^n-1\right)&>\gcd\left(2^m-1,2^n-1\right) \hspace{12mm} &&\text{by }\textbf{Claim 1} \\ \gcd\left(2^{2m}-1,2^n-1\right)&>\gcd\left(2^m-1,2^n-1\right) \\ 2^{\gcd(2m,n)}-1&>2^{\gcd(m,n)}-1 &&\text{by }\textbf{Claim 2} \\ \gcd(2m,n)&>\gcd(m,n), \end{align*}
which implies that $n$ must have more factors of $2$ than $m$ does.

 We construct the following table for the first $30$ positive integers:
 \[\begin{array}{c|c|c} && \\ [-2.5ex] \boldsymbol{\#}\textbf{ of Factors of }\boldsymbol{2} & \textbf{Numbers} & \textbf{Count} \\ \hline && \\ [-2.25ex]  0 & 1,3,5,7,9,11,13,15,17,19,21,23,25,27,29 & 15 \\  && \\ [-2.25ex]  1 & 2,6,10,14,18,22,26,30 & 8 \\ && \\ [-2.25ex]  2 & 4,12,20,28 & 4 \\    && \\ [-2.25ex]  3 & 8,24 & 2 \\  && \\ [-2.25ex]  4 & 16 & 1 \\ \end{array}\]To count the ordered pairs $(m,n),$ we perform casework on the number of factors of $2$ that $m$ has:

 <ol style="margin-left: 1.5em;"> \item If $m$ has $0$ factors of $2,$ then $m$ has $15$ options and $n$ has $8+4+2+1=15$ options. So, this case has $15\cdot15=225$ ordered pairs. 

 \item If $m$ has $1$ factor of $2,$ then $m$ has $8$ options and $n$ has $4+2+1=7$ options. So, this case has $8\cdot7=56$ ordered pairs. 

 \item If $m$ has $2$ factors of $2,$ then $m$ has $4$ options and $n$ has $2+1=3$ options. So, this case has $4\cdot3=12$ ordered pairs. 

 \item If $m$ has $3$ factors of $2,$ then $m$ has $2$ options and $n$ has $1$ option. So, this case has $2\cdot1=2$ ordered pairs.</ol>Together, the answer is $225+56+12+2=\boxed{295}.$

 





\textbf{Solution 2}

Consider any ordered pair $(m,n)$ such that $\gcd(2^m+1, 2^n-1) > 1$. There must exist some odd number $p\ne 1$ such that $2^m \equiv -1 \pmod{p}$ and $2^n \equiv 1 \pmod{p}$. Let $d$ be the order of $2$ modulo $p$. Note that $2^{2m} \equiv 1 \pmod{p}$. From this, we can say that $2m$ and $n$ are both multiples of $d$, but $m$ is not. Thus, we have $v_2(n) \ge v_2(d)$ and $v_2(m) + 1 = v_2(d)$.  Substituting the latter equation into the inequality before gives $v_2(n) \ge v_2(m)+1$. Since $v_2(n)$ and $v_2(m)$ are integers, this implies $v_2(n)>v_2(m)$. The rest of the solution now proceeds as in Solution 1.

 \textbf{Remarks} \textbf{Claim 1 (GCD Property)} \textbf{If $\boldsymbol{r,s,}$ and $\boldsymbol{t}$ are positive integers such that $\boldsymbol{\gcd(r,s)=1,}$ then $\boldsymbol{\gcd(r,t)\cdot\gcd(s,t)=\gcd(rs,t).}$}

 As $r$ and $s$ are relatively prime (have no prime divisors in common), this property is intuitive.

 \textbf{Claim 2 (Olympiad Number Theory Lemma)} \textbf{If $\boldsymbol{u,a,}$ and $\boldsymbol{b}$ are positive integers such that $\boldsymbol{u\geq2,}$ then $\boldsymbol{\gcd\left(u^a-1,u^b-1\right)=u^{\gcd(a,b)}-1.}$}

 There are two proofs to this claim, as shown below.

 \textbf{Claim 2 Proof 1 (Euclidean Algorithm)} If $a=b,$ then $\gcd(a,b)=a=b,$ from which the claim is clearly true.

 Otherwise, let $a>b$ without the loss of generality. For all integers $p$ and $q$ such that $p>q>0,$ the Euclidean Algorithm states that 
 \[\gcd(p,q)=\gcd(p-q,q)=\cdots=\gcd(p\operatorname{mod}q,q).\] We apply this result repeatedly to reduce the larger number: 
 \[\gcd\left(u^a-1,u^b-1\right)=\gcd\left(u^a-1-u^{a-b}\left(u^b-1\right),u^b-1\right)=\gcd\left(u^{a-b}-1,u^b-1\right).\]Continuing, we have
 \begin{align*} \gcd\left(u^a-1,u^b-1\right)&=\gcd\left(u^{a-b}-1,u^b-1\right) \\ & \ \vdots \\ &=\gcd\left(u^{\gcd(a,b)}-1,u^{\gcd(a,b)}-1\right) \\ &=u^{\gcd(a,b)}-1, \end{align*}
from which the proof is complete.

 \textbf{Claim 2 Proof 2 (Bézout's Identity)} Let $d=\gcd\left(u^a-1,u^b-1\right).$ It follows that $u^a\equiv1\pmod{d}$ and $u^b\equiv1\pmod{d}.$

 By Bézout's Identity, there exist integers $x$ and $y$ such that $ax+by=\gcd(a,b),$ so 
 \[u^{\gcd(a,b)}=u^{ax+by}=(u^a)^x\cdot(u^b)^y\equiv1\pmod{d},\]from which $u^{\gcd(a,b)}-1\equiv0\pmod{d}.$ We know that $u^{\gcd(a,b)}-1\geq d.$

 Next, we notice that
 \begin{align*} u^a-1&=\left(u^{\gcd(a,b)}-1\right)\left(u^{a-\gcd{(a,b)}}+u^{a-2\gcd{(a,b)}}+u^{a-3\gcd{(a,b)}}+\cdots+1\right), \\ u^b-1&=\left(u^{\gcd(a,b)}-1\right)\left(u^{b-\gcd{(a,b)}}+u^{b-2\gcd{(a,b)}}+u^{b-3\gcd{(a,b)}}+\cdots+1\right). \end{align*}
Since $u^{\gcd(a,b)}-1$ is a common divisor of $u^a-1$ and $u^b-1,$ we conclude that $u^{\gcd(a,b)}-1=d,$ from which the proof is complete.

 




	\item (\textit{AIME II 2021}) Two spheres with radii $36$ and one sphere with radius $13$ are each externally tangent to the other two spheres and to two different planes $\mathcal{P}$ and $\mathcal{Q}$. The intersection of planes $\mathcal{P}$ and $\mathcal{Q}$ is the line $\ell$. The distance from line $\ell$ to the point where the sphere with radius $13$ is tangent to plane $\mathcal{P}$ is $\frac{m}{n}$, where $m$ and $n$ are relatively prime positive integers. Find $m + n$.


 
 \textbf{Answer: }335



\textbf{Solution 1 (Similar Triangles and Pythagorean Theorem)}

This solution refers to the \textbf{Diagram} section.

 As shown below, let $O_1,O_2,O_3$ be the centers of the spheres (where sphere $O_3$ has radius $13$) and $T_1,T_2,T_3$ be their respective points of tangency to plane $\mathcal{P}.$ Let $\mathcal{R}$ be the plane that is determined by $O_1,O_2,$ and $O_3.$ Suppose $A$ is the foot of the perpendicular from $O_3$ to line $\ell,$ so $\overleftrightarrow{O_3A}$ is the perpendicular bisector of $\overline{O_1O_2}.$ We wish to find $T_3A.$
                    \begin{figure}[H]
                        \centering
                        \includegraphics[width=0.5\textwidth]{Solution/AIME_2021_II_10_Solution_Figure_1.png}
                    \end{figure}
                    Note that:

 <ol style="margin-left: 1.5em;"> \item In $\triangle O_1O_2O_3,$ we get $O_1O_2=72$ and $O_1O_3=O_2O_3=49.$ 

 \item Both $\triangle O_1O_2O_3$ and $\overline{O_3A}$ lie in plane $\mathcal{R}.$ Both $\triangle T_1T_2T_3$ and $\overline{T_3A}$ lie in plane $\mathcal{P}.$ 

 \item By symmetry, since planes $\mathcal{P}$ and $\mathcal{Q}$ are reflections of each other about plane $\mathcal{R},$ the three planes are concurrent to line $\ell.$ 

 \item Since $\overline{O_1T_1}\perp\mathcal{P}$ and $\overline{O_3T_3}\perp\mathcal{P},$ it follows that $\overline{O_1T_1}\parallel\overline{O_3T_3},$ from which $O_1,O_3,T_1,$ and $T_3$ are coplanar.

 </ol>Now, we focus on cross-sections $O_1O_3T_3T_1$ and $\mathcal{R}:$

 <ol style="margin-left: 1.5em;"> <li><i>\textbf{In the three-dimensional space, the intersection of a line and a plane must be exactly one of the empty set, a point, or a line.}</i>Clearly, cross-section $O_1O_3T_3T_1$ intersects line $\ell$ at exactly one point. Furthermore, as the intersection of planes $\mathcal{R}$ and $\mathcal{P}$ is line $\ell,$ we conclude that $\overrightarrow{O_1O_3}$ and $\overrightarrow{T_1T_3}$ must intersect line $\ell$ at the same point. Let $B$ be the point of concurrency of $\overrightarrow{O_1O_3},\overrightarrow{T_1T_3},$ and line $\ell.$

 </li> \item In cross-section $\mathcal{R},$ let $C$ be the foot of the perpendicular from $O_1$ to line $\ell,$ and $D$ be the foot of the perpendicular from $O_3$ to $\overline{O_1C}.$

 </ol>We have the following diagram:
                    \begin{figure}[H]
                        \centering
                        \includegraphics[width=0.5\textwidth]{Solution/AIME_2021_II_10_Solution_Figure_2.png}
                    \end{figure}
                    In cross-section $O_1O_3T_3T_1,$ since $\overline{O_1T_1}\parallel\overline{O_3T_3}$ as discussed, we obtain $\triangle O_1T_1B\sim\triangle O_3T_3B$ by AA, with the ratio of similitude $\frac{O_1T_1}{O_3T_3}=\frac{36}{13}.$ Therefore, we get $\frac{O_1B}{O_3B}=\frac{49+O_3B}{O_3B}=\frac{36}{13},$ or $O_3B=\frac{637}{23}.$

 In cross-section $\mathcal{R},$ note that $O_1O_3=49$ and $DO_3=\frac{O_1O_2}{2}=36.$ Applying the Pythagorean Theorem to right $\triangle O_1DO_3,$ we have $O_1D=\sqrt{1105}.$ Moreover, since $\ell\perp\overline{O_1C}$ and $\overline{DO_3}\perp\overline{O_1C},$ we obtain $\ell\parallel\overline{DO_3}$ so that $\triangle O_1CB\sim\triangle O_1DO_3$ by AA, with the ratio of similitude $\frac{O_1B}{O_1O_3}=\frac{49+\frac{637}{23}}{49}.$ Therefore, we get $\frac{O_1C}{O_1D}=\frac{\sqrt{1105}+DC}{\sqrt{1105}}=\frac{49+\frac{637}{23}}{49},$ or $DC=\frac{13\sqrt{1105}}{23}.$

 Finally, note that $\overline{O_3T_3}\perp\overline{T_3A}$ and $O_3T_3=13.$ Since quadrilateral $DCAO_3$ is a rectangle, we have $O_3A=DC=\frac{13\sqrt{1105}}{23}.$ Applying the Pythagorean Theorem to right $\triangle O_3T_3A$ gives $T_3A=\frac{312}{23},$ from which the answer is $312+23=\boxed{335}.$

 





\textbf{Solution 2 (Pythagorean Theorem)}

The centers of the three spheres form a $49$-$49$-$72$ triangle. Consider the points at which the plane is tangent to the two bigger spheres; the line segment connecting these two points should be parallel to the $72$ side of this triangle. Take its midpoint $M$, which is $36$ away from the midpoint $A$ of the $72$ side, and connect these two midpoints.

 Now consider the point at which the plane is tangent to the small sphere, and connect $M$ with the small sphere's tangent point $B$. Extend $\overline{MB}$ through $B$ until it hits the ray from $A$ through the center of the small sphere (convince yourself that these two intersect). Call this intersection $D$, the center of the small sphere $C$, we want to find $BD$.

 By Pythagoras, $AC=\sqrt{49^2-36^2}=\sqrt{1105}$, and we know that $MA=36$ and $BC=13$. We know that $\overline{MA}$ and $\overline{BC}$ must be parallel, using ratios we realize that $CD=\frac{13}{23}\sqrt{1105}$. Apply the Pythagorean theorem to $\triangle BCD$, $BD=\frac{312}{23}$, so $312 + 23 = \boxed{335}$.

 





\textbf{Solution 3 (Proportion)}

This solution refers to the \textbf{Diagram} section.
                    \begin{figure}[H]
                        \centering
                        \includegraphics[width=0.5\textwidth]{Solution/AIME_2021_II_10_Solution_Figure_3.png}
                    \end{figure}
                    The isosceles triangle of centers $O_1 O_2 O$ ($O$ is the center of sphere of radii $13$) has sides  $O_1 O = O_2 O = 36 + 13 = 49,$ and  $O_1 O_2 = 36 + 36 = 72.$ 

 Let $N$ be the midpoint  $O_1 O_2$.

 The isosceles triangle of points of tangency $T_1 T_2 T$  has sides $T_1 T = T_2 T = 2 \sqrt{13 \cdot 36} = 12 \sqrt{13}$ and $T_1 T_2 = 72.$ 

 Let $M$ be the midpoint $T_1 T_2.$

 The height $TM$  is $\sqrt {12^2 \cdot 13 - 36^2} = 12 \sqrt {13-9} = 24.$

 The tangents of the half-angle between the planes is $\frac {TO}{AT} = \frac {MN - TO}{TM},$ so $\frac {13}{AT} = \frac {36 - 13}{24},$ 
 \[AT = \frac{24\cdot 13}{23} = \frac {312}{23} \implies  312 + 23 = \boxed{335}.\]\textbf{vladimir.shelomovskii@gmail.com, vvsss}

 





\textbf{Solution 4 (Coordinate Geometry)}

Let the centers of the spheres be $O_1,O_2,O_3$. We establish a 3D coordinate system with $\mathcal{P}$ being the $x-y$ plane, and here is our setup:\begin{align*}O_1&=(0,-36,36)\\O_2&=(0,36,36)\\O_3&=(x_1,0,13)\end{align*}
By the distance formula, 
 \[x_{1}^{2}+36^2+23^2=49^2\implies x_{1}=24.\]Let the equation of $\mathcal{Q}$ be $ax+by+cz+d=0$. WLOG, we can assume $\sqrt{a^{2}+b^{2}+c^2}=1$. Applying the distance formula on the $3$ centers yields the system\begin{align*}\lvert -36b+36c+d\rvert=36\\\lvert 36b+36c+d\rvert=36\\\lvert 24a+13c+d\rvert=13\end{align*}
By symmetry, $\mathcal{Q}$ is parallel to the $y$-axis, hence $b=0$. Furthermore, since all spheres lie entirely on one side of $\mathcal{Q}$, we can choose the direction of the normal vector such that the inner expressions are positive. Dropping the absolute values gives\begin{align*}36c+d=36\\24a+13c+d=13\\a^{2}+c^{2}=1\end{align*}
Solving yields $a=-\frac{1104}{1105},c=-\frac{47}{1105},d=\frac{41472}{1105}$.To find the intersection line $\ell$ between $\mathcal{Q}$ and $\mathcal{P}$, we set $z=0$, which gives the equation $x=-\frac{d}{a}=\frac{864}{23}$. The small sphere touches $\mathcal{P}$ at $(24,0,0)$. The distance from this point of tangency to $\ell$ is thus 
 \[\frac{864}{23}-24=\frac{312}{23}\]The requested sum is $312+23=\boxed{335}$.

 




	\item (\textit{AIME II 2021}) A teacher was leading a class of four perfectly logical students. The teacher chose a set $S$ of four integers and gave a different number in $S$ to each student. Then the teacher announced to the class that the numbers in $S$ were four consecutive two-digit positive integers, that some number in $S$ was divisible by $6$, and a different number in $S$ was divisible by $7$. The teacher then asked if any of the students could deduce what $S$ is, but in unison, all of the students replied no.


 However, upon hearing that all four students replied no, each student was able to determine the elements of $S$. Find the sum of all possible values of the greatest element of $S$.


 
 \textbf{Answer: }258



\textbf{Solution 1}

Note that $\operatorname{lcm}(6,7)=42.$ It is clear that $42\not\in S$ and $84\not\in S,$ otherwise the three other elements in $S$ are divisible by neither $6$ nor $7.$

 In the table below, the multiples of $6$ are colored in yellow, and the multiples of $7$ are colored in green. By the least common multiple, we obtain cycles: If $n$ is a possible maximum value of $S,$ then $n+42$ must be another possible maximum value of $S,$ and vice versa. By observations, we circle all possible maximum values of $S.$
                    \begin{figure}[H]
                        \centering
                        \includegraphics[width=0.5\textwidth]{Solution/AIME_2021_II_11_Solution_Figure_1.png}
                    \end{figure}
                    From the second row of the table above, we perform casework on the possible maximum value of $S:$
 \[\begin{array}{c||c|c|l} & & & \\ [-2.5ex] \textbf{Max Value} & \boldsymbol{S} & \textbf{Valid?} & \hspace{16.25mm}\textbf{Reasoning/Conclusion} \\ [0.5ex] \hline & & & \\ [-2ex] 49 & \{46,47,48,49\} & & \text{The student who gets } 46 \text{ will reply yes.} \\ 50 & \{47,48,49,50\} & \checkmark & \text{Another possibility is } S=\{89,90,91,92\}. \\ 51 & \{48,49,50,51\} & & \text{The student who gets } 51 \text{ will reply yes.} \\ 56 & \{53,54,55,56\} & & \text{The student who gets } 53 \text{ will reply yes.} \\ 57 & \{54,55,56,57\} & & \text{The student who gets } 57 \text{ will reply yes.} \\ 63 & \{60,61,62,63\} & & \text{The students who get } 60,61,62 \text{ will reply yes.} \\ 66 & \{63,64,65,66\} & & \text{The students who get } 64,65,66 \text{ will reply yes.} \\ 72 & \{69,70,71,72\} & & \text{The student who gets } 69 \text{ will reply yes.} \\ 73 & \{70,71,72,73\} & & \text{The student who gets } 73 \text{ will reply yes.} \\ 78 & \{75,76,77,78\} & & \text{The student who gets } 75 \text{ will reply yes.} \\ 79 & \{76,77,78,79\} & \checkmark & \text{Another possibility is } S=\{34,35,36,37\}. \\ 80 & \{77,78,79,80\} & & \text{The student who gets } 80 \text{ will reply yes.} \end{array}\]Finally, all possibilities for $S$ are $\{34,35,36,37\}, \{47,48,49,50\}, \{76,77,78,79\},$ and $\{89,90,91,92\},$ from which the answer is $37+50+79+92=\boxed{258}.$

 <u>\textbf{Remarks}</u>

 <ol style="margin-left: 1.5em;"> \item Alternatively, we can reconstruct the second table in this solution as follows, where Y and N denote the replies of "yes" and "no", respectively. Notice that this table has some kind of symmetry!
                    \begin{figure}[H]
                        \centering
                        \includegraphics[width=0.5\textwidth]{Solution/AIME_2021_II_11_Solution_Figure_2.png}
                    \end{figure}
                     \item As a confirmation, we can verify that each student will be able to deduce what $S$ is upon hearing the four replies of "no" in unison. For example, if $S=\{47,48,49,50\},$ then all students will know that no one gets $46$ or $51,$ otherwise that student will reply yes (as discussed). Therefore, all students will conclude that $S$ has only one possibility.

 </ol>





\textbf{Solution 2}

We know right away that $42\not\in S$ and $84\not\in S$ as stated in Solution 1.

 To get a feel for the problem, let’s write out some possible values of $S$ based on the teacher’s remarks. The first multiple of 7 that is two-digit is 14. The closest multiple of six from 14 is 12, and therefore there are two possible sets of four consecutive integers containing 12 and 14; $\{11,12,13,14\}$ and $\{12,13,14,15\}$. Here we get our first crucial idea; that if the multiples of 6 and 7 differ by two, there will be 2 possible sets of $S$ without any student input. Similarly, if they differ by 3, there will be only 1 possible set, and if they differ by 1, 3 possible sets.

 Now we read the student input. Each student says they can’t figure out what $S$ is just based on the teacher’s information, which means each student has to have a number that would be in 2 or 3 of the possible sets (This is based off of the first line of student input). However, now that each student knows that all of them have numbers that fit into more than one possible set, this means that S cannot have two possible sets because otherwise, when shifting from one set to the other, one of the end numbers would not be in the shifted set, but we know each number has to fall in two or more possible sets. For example, take $\{11,12,13,14\}$ and $\{12,13,14,15\}$. The numbers at the end, 11 and 15, only fall in one set, but each number must fall in at least two sets. This means that there must be three possible sets of S, in which case the actual S would be the middle S.Take for example $\{33,34,35,36\}$, $\{34,35,36,37\}$, and $\{35,36,37,38\}$. 37 and 34 fall in two sets while 35 and 36 fall in all three sets, so the condition is met. Now, this means that the multiple of 6 and 7 must differ by 1. Since 42 means the difference is 0, when you add/subtract 6 and 7, you will obtain the desired difference of 1, and similarly adding/subtracting 6 or 7 from 84 will also obtain the difference of 1. Thus there are four possible sets of $S$; $\{34,35,36,37\}$, $\{47,48,49,50\}$, $\{76,77,78,79\}$ and $\{89,90,91,92.\}$. Therefore the sum of the greatest elements of the possible sets $S$ is $37+50+79+92=\boxed{258}$

 





\textbf{Solution 3}

In a solution that satisfies these constraints, the multiple of 6 must be adjacent to multiple of 7. The other two numbers must be on either side. 

 WLOG assume the set is $\{a,6j,7k,b\}$. The student with numbers $a$, $6j$, and $7k$ can think the set is $\{a-1, a,6j,7k\}$ or $\{a,6j,7k,b\}$, and the students with number $6j$, $7k$, and $b$ can think the set is $\{a,6j,7k,b\}$ or $\{6j,7k,b, b+1\}$. Therefore, none of the students know the set for sure. 

 Playing around with the arrangement of the multiple of 6 and multiple of 7 shows that this is the only configuration viewed as ambiguous to all the students. (Therefore when they hear nobody else knows either, they can find out it is this configuration)

 Considering $S$ as $\{a,6j,7k,b\}$,b is 2 mod 6 and 1 mod 7, so $b$ is 8 mod 42. (since it is all 2-digit, the values are either 50 or 92).

 Similarly, considering $S$ as $\{a,7j,6k,b\}$, $b$ is 1 mod 6 and 2 mod 7, so $b$ is 37 mod 42. The values that satisfy that are 37 and 79. 

 The total sum of all these values is therefore $50+92+37+79=\boxed{258}$.

 This solution will even work when the bounds (in this question, 2-digit so <100) are much larger and it is impractical to perform casework.

 





\textbf{Solution 4 (Logic)}

Consider the tuple $(a, a+1, a+2, a+3)$ as a possible $S$. If one of the values in $S$ is $3$ or $4 \pmod{7}$, observe the student will be able to deduce $S$ with no additional information. This is because, if a value is $b = 3 \pmod{7}$ and $S$ contains a $0 \pmod{7}$, then the values of $S$ must be $(b-3, b-2, b-1, b)$. Similarly, if we are given a $b \equiv 4 \pmod{7}$ and we know that $0 \pmod{7}$ is in $S$, $S$ must be $(b, b+1, b+2, b+3).$ Hence, the only possibility for $a$ is $5, 6 \pmod{7}.$

 In either case, we are guaranteed there is a $6, 0, 1 \pmod{7}$ value in $7$. The difference comes down to if there is a $5 \pmod{7}$ value or a $1 \pmod{7}$ value. The person receiving such value will be able to determine all of $S$ but the $6, 0, 1 \pmod{7}$ people will not be able to differentiate the two cases ... yet.

 Now consider which value among the consecutive integers is $c \equiv 3 \pmod{6}$, if any. The person will know that $S$ is either $(c-3, c-2, c-1, c)$ or $(c, c+1, c+2, c+3)$ to have a $0 \pmod{6}$ value in $S$. Neither the $6, 1 \pmod{7}$ person can be $3 \pmod{7}$, else they can decipher what $S$ is right off the bat by considering which set has $0 \pmod{7}.$ This translates to the possible $5 \pmod{7}$ or $1 \pmod{7}$ person cannot be $0 \pmod{6}$. We are given that $0 \pmod{7}$ and $0 \pmod{6}$ cannot be the same person. Hence we conclude one of the $6 \pmod{7}$ or $1 \pmod{7}$ must be the $0 \pmod{6}$ person.

 Let the $6 \pmod{7}$ person be $0 \pmod{6}.$ Then--hypothetical--$c \equiv 2 \pmod{7}$ person is $3 \pmod{6}.$ After the first round, the $6, 0, 1 \pmod{7}$ people realize that $2 \pmod{7}$ is not in $S$ else they would have deduced $S$ by noting $S$ was either $(c-3, c-2, c-1, c)$ or $(c, c+1, c+2, c+3)$ to have a $0 \pmod{6}$ and choosing the former based on where $0 \pmod{7}$ is. Hence they figure out $S$ by knowing $5 \pmod{7}.$ So $a \equiv 5 \pmod{7}$ and $a \equiv 5 \pmod{6}$ (from $6 \pmod{7}$ person being $0 \pmod{6}$).

 Similarly, if  $1 \pmod{7}$ person is $0 \pmod{6}$ we find that $a \equiv 6 \pmod{7}$ (so a $2 \pmod{7}$ is in $S$) and $a \equiv 4 \pmod{6}.$

 By CRT, the possibilities are $a \equiv 5,  34 \pmod{42}.$ The sum of the greatest values of $S$ are the sum of $a + 3$ and so we get $(34 + 3) + (47 + 3) + (76 + 3) + (89 + 3) = \boxed{258}.$ 

 





\textbf{Solution 5 (THINK)}

Start by writing down every two-digit multiple of $6$ and $7$. First, $42$ and $84$ will not be in $S$ because we need two distinct numbers where one is a multiple of $6$ and ANOTHER is a multiple of $7$. Furthermore, it is clear at this point that if the multiples of $6$ and $7$ are exactly $3$ apart, then it will be obvious what $S$ is; it is just the smaller multiple to the larger multiple. Therefore, we cross out numbers $18, 21, 24, 60, 63, 66$. 

 That leaves us with pairs of multiples of $6$ and $7$ that are exactly $1$ or $2$ apart. In this case, if the multiples of $6$ and $7$ are two apart, such as $12$ and $14$, then the non-multiples will be able to deduce what the set is; this leaves that the $6$ and $7$ must be one apart.

 Specifically, these are pairs $(35, 36), (48, 49), (77, 78),$ and $(90, 91)$. In each of these pairs, respectively, the largest number will be one greater than the larger multiple, namely $37, 50, 79, 92$. Therefore, our answer is $37 + 50 + 79 + 92 = \boxed{258}$. 

 \textbf{Animated Solution by Interstigation} 




	\item (\textit{AIME II 2021}) A convex quadrilateral has area $30$ and side lengths $5, 6, 9,$ and $7,$ in that order. Denote by $\theta$ the measure of the acute angle formed by the diagonals of the quadrilateral. Then $\tan \theta$ can be written in the form $\frac{m}{n}$, where $m$ and $n$ are relatively prime positive integers. Find $m + n$.


 
 \textbf{Answer: }047



\textbf{Solution 1 (Sines and Cosines)}

Since we are asked to find $\tan \theta$, we can find $\sin \theta$ and $\cos \theta$ separately and use their values to get $\tan \theta$. We can start by drawing a diagram. Let the vertices of the quadrilateral be $A$, $B$, $C$, and $D$. Let $AB = 5$, $BC = 6$, $CD = 9$, and $DA = 7$. Let $AX = a$, $BX = b$, $CX = c$, and $DX = d$. We know that $\theta$ is the acute angle formed between the intersection of the diagonals $AC$ and $BD$.
                    \begin{figure}[H]
                        \centering
                        \includegraphics[width=0.5\textwidth]{Solution/AIME_2021_II_12_Solution_Figure_1.png}
                    \end{figure}
                    We are given that the area of quadrilateral $ABCD$ is $30$. We can express this area using the areas of triangles $AXB$, $BXC$, $CXD$, and $DXA$. Since we want to find $\sin \theta$ and $\cos \theta$, we can represent these areas using $\sin \theta$ as follows:
 \begin{align*} 30 &=[ABCD] \\ &=[AXB] + [BXC] + [CXD] + [DXA] \\ &=\frac{1}{2} ab \sin (\angle AXB) + \frac{1}{2} bc \sin (\angle BXC) + \frac{1}{2} cd \sin (\angle CXD) + \frac{1}{2} da \sin (\angle AXD) \\ &=\frac{1}{2} ab \sin (180^\circ - \theta) + \frac{1}{2} bc \sin (\theta) + \frac{1}{2} cd \sin (180^\circ - \theta) + \frac{1}{2} da \sin (\theta). \end{align*}
We know that $\sin (180^\circ - \theta) = \sin \theta$. Therefore it follows that:
 \begin{align*} 30 &=\frac{1}{2} ab \sin (180^\circ - \theta) + \frac{1}{2} bc \sin (\theta) + \frac{1}{2} cd \sin (180^\circ - \theta) + \frac{1}{2} da \sin (\theta) \\ &=\frac{1}{2} ab \sin (\theta) + \frac{1}{2} bc \sin (\theta) + \frac{1}{2} cd \sin (\theta) + \frac{1}{2} da \sin (\theta) \\ &=\frac{1}{2}\sin\theta (ab + bc + cd + da). \end{align*}
From here we see that $\sin \theta = \frac{60}{ab + bc + cd + da}$. Now we need to find $\cos \theta$. Using the Law of Cosines on each of the four smaller triangles, we get following equations:
 \begin{align*}   5^2 &= a^2 + b^2 - 2ab\cos(180^\circ-\theta), \\   6^2 &= b^2 + c^2 - 2bc\cos \theta, \\   9^2 &= c^2 + d^2 - 2cd\cos(180^\circ-\theta), \\   7^2 &= d^2 + a^2 - 2da\cos \theta.   \end{align*}
We know that $\cos (180^\circ - \theta) = -\cos \theta$ for all $\theta$. We can substitute this value into our equations to get:
 \begin{align*}   5^2 &= a^2 + b^2 + 2ab\cos \theta, &&(1) \\   6^2 &= b^2 + c^2 - 2bc\cos \theta, &&(2) \\   9^2 &= c^2 + d^2 + 2cd\cos \theta, &&(3) \\   7^2 &= d^2 + a^2 - 2da\cos \theta. &&(4) \end{align*}
If we subtract $(2)+(4)$ from $(1)+(3)$, the squared terms cancel, leaving us with:
 \begin{align*} 5^2 + 9^2 - 6^2 - 7^2 &= 2ab \cos \theta + 2bc \cos \theta + 2cd \cos \theta + 2da \cos \theta \\  21 &= 2\cos \theta (ab  + bc + cd + da). \end{align*}
From here we see that $\cos \theta = \frac{21/2}{ab + bc + cd + da}$.

 Since we have figured out $\sin \theta$ and $\cos \theta$, we can calculate $\tan \theta$: 
 \[\tan \theta = \frac{\sin \theta}{\cos \theta} = \frac{\frac{60}{ab + bc + cd + da}}{\frac{21/2}{ab + bc + cd + da}} = \frac{60}{21/2} = \frac{120}{21} = \frac{40}{7}.\]Therefore our answer is $40 + 7 = \boxed{047}$.

 





\textbf{Solution 2 (Right Triangles)}

In convex quadrilateral $ABCD,$ let $AB=5,BC=6,CD=9,$ and $DA=7.$ Let $A'$ and $C'$ be the feet of the perpendiculars from $A$ and $C,$ respectively, to $\overline{BD}.$ We obtain the following diagram:
                    \begin{figure}[H]
                        \centering
                        \includegraphics[width=0.5\textwidth]{Solution/AIME_2021_II_12_Solution_Figure_2.png}
                    \end{figure}
                    Let $BC'=p,C'E=q,EA'=r,A'D=s,AA'=h_1,$ and $CC'=h_2.$ We apply the Pythagorean Theorem to right triangles $\triangle ABA',\triangle BCC',\triangle CDC',$ and $\triangle DAA',$ respectively:
 \[\begin{array}{ccccccccccccccccc} (p+q+r)^2&+&h_1^2&=&5^2, &&&&&&&&&&&&\hspace{36mm}(1) \\ [1ex] p^2&+&h_2^2&=&6^2, &&&&&&&&&&&&\hspace{36mm}(2) \\ [1ex] (q+r+s)^2&+&h_2^2&=&9^2, &&&&&&&&&&&&\hspace{36mm}(3) \\ [1ex] s^2&+&h_1^2&=&7^2. &&&&&&&&&&&&\hspace{36mm}(4) \end{array}\]Let the brackets denote areas. We get
 \begin{align*} [ABD]+[CBD]&=[ABCD] \\ \frac12(p+q+r+s)h_1+\frac12(p+q+r+s)h_2&=30 \\ \frac12(p+q+r+s)(h_1+h_2)&=30 \\ (p+q+r+s)(h_1+h_2)&=60. \hspace{49.25mm}(5) \end{align*}
We subtract $(2)+(4)$ from $(1)+(3):$
 \begin{align*} (p+q+r)^2+(q+r+s)^2-p^2-s^2&=21 \\ \left[(p+q+r)^2-s^2\right]+\left[(q+r+s)^2-p^2\right]&=21 \\ (p+q+r+s)(p+q+r-s)+(p+q+r+s)(-p+q+r+s)&=21 \\ (p+q+r+s)(2q+2r)&=21 \\ 2(p+q+r+s)(q+r)&=21 \\ (p+q+r+s)(q+r)&=\frac{21}{2}. \hspace{9.5mm}(6) \end{align*}
From right triangles $\triangle AEA'$ and $\triangle CEC',$ we have $\tan\theta=\frac{h_1}{r}=\frac{h_2}{q}.$ It follows that
 \begin{alignat*}{8} \tan\theta&=\frac{h\_1}{r}\qquad&\implies\qquad h\_1&=r\tan\theta, \hspace{64mm}&(1\star)\\ \tan\theta&=\frac{h\_2}{q}\qquad&\implies\qquad h\_2&=q\tan\theta. &(2\star) \end{alignat*}Finally, we divide $(5)$ by $(6):$
 \begin{align*} \frac{h_1+h_2}{q+r}&=\frac{40}{7} \\ \frac{r\tan\theta+q\tan\theta}{q+r}&=\frac{40}{7} \hspace{15mm} &&\text{by }(1\star)\text{ and }(2\star)\\ \frac{(r+q)\tan\theta}{q+r}&=\frac{40}{7} \\ \tan\theta&=\frac{40}{7}, \end{align*}
from which the answer is $40+7=\boxed{047}.$

 





\textbf{Solution 3 (Bretschneider's Formula)}

\textbf{Bretschneider's Formula} 
                    \begin{figure}[H]
                        \centering
                        \includegraphics[width=0.5\textwidth]{Solution/AIME_2021_II_12_Solution_Figure_3.png}
                    \end{figure}
                    Given quadrilateral $ABCD$, let, $a, b, c, d$, be the sides, $s$ the semiperimeter, and $u, v$, the diagonals. Then the area, $K$, is given by 
 \[K = \frac 14 \sqrt {4u^2v^2-(b^2+d^2-a^2-c^2)^2}\]

 





\textbf{Solution}

By Bretschneider's Formula, 
 \[30=\frac{1}{4}\sqrt{4u^2v^2-(b^2+d^2-a^2-c^2)^2}=\frac{1}{4}\sqrt{4u^2v^2-441}.\] Thus, $uv=\frac{3\sqrt{1649}}{2}$. Also, 
 \[[ABCD]=\frac 12 \cdot uv\sin{\theta};\] solving for $\sin{\theta}$ yields $\sin{\theta}=\frac{40}{\sqrt{1649}}$. Since $\theta$ is acute, $\cos{\theta}$ is positive, from which $\cos{\theta}=\frac{7}{\sqrt{1649}}$. Solving for $\tan{\theta}$ yields 
 \[\tan{\theta}=\frac{\sin{\theta}}{\cos{\theta}}=\frac{40}{7},\] for a final answer of $\boxed{047}$. 

 





\textbf{Solution 4 (Symmetry)}


                    \begin{figure}[H]
                        \centering
                        \includegraphics[width=0.5\textwidth]{Solution/AIME_2021_II_12_Solution_Figure_4.png}
                    \end{figure}
                    <i>\textbf{Claim}</i>

 Given an inscribed quadrilateral $ABCD$ with sides $AB = a, BC = b, CD = c,$ and $DA = d.$ Prove that the $\angle \theta < 90^\circ$ between the diagonals is given by
 \begin{align*}2(ac + bd) \cos \theta =  {|d^2 – c^2 + b^2 – a^2|}.\end{align*}
<i>\textbf{Proof}</i>

 Let the point $B'$ be symmetric to $B$ with respect to the perpendicular bisector $AC.$ Then the quadrilateral $AB'CD$ is an inscribed one, $AB' = b, B'C = a.$

 
 \[2 \angle AEB = \overset{\Large\frown} {AB}  + \overset{\Large\frown} {CD}.\]
 \begin{align*} 2\angle B'AD = \overset{\Large\frown} {B'C} + \overset{\Large\frown} {CD} = \overset{\Large\frown} {AB} + \overset{\Large\frown} {CD} \implies \angle AEB = \angle B'AD.\end{align*}


 We apply the Law of Cosines to $\triangle AB'D$ and $\triangle CB'D$:
 \begin{align*} B'D^2 = AD^2 + AB'^2 – 2 AD \cdot AB' \cos \theta, \end{align*}

 \begin{align*} B'D^2 = CD^2 + CB'^2 + 2 CD \cdot CB' \cos \theta,\end{align*}

 \begin{align*} d^2 + b^2 – 2 bd \cos \theta = c^2 + a^2 + 2ac \cos \theta,\end{align*}

 \begin{align*} 2(ac + bd) \cos \theta = |d^2 – c^2 + b^2 – a^2|.\end{align*}


 \textbf{vladimir.shelomovskii@gmail.com, vvsss}

 \textbf{Note 1} By generalization, the tangent of the acute angle formed by the diagonals is 
 \[\left|\frac{4A}{a^2-b^2+c^2-d^2}\right|.\]

 




	\item (\textit{AIME II 2021}) Find the least positive integer $n$ for which $2^n + 5^n - n$ is a multiple of $1000$.


 
 \textbf{Answer: }797



\textbf{Solution 1}

Recall that $1000$ divides this expression if $8$ and $125$ both divide it. It should be fairly obvious that $n \geq 3$; so we may break up the initial condition into two sub-conditions.

 (1) $5^n \equiv n \pmod{8}$. Notice that the square of any odd integer is $1$ modulo $8$ (proof by plugging in $1^2,3^2,5^2,7^2$ into modulo $8$), so the LHS of this expression goes $5,1,5,1,\ldots$, while the RHS goes $1,2,3,4,5,6,7,8,1,\ldots$. The cycle length of the LHS is $2$, RHS is $8$, so the cycle length of the solution is $\operatorname{lcm}(2,8)=8$. Indeed, the $n$ that solve this congruence are exactly those such that $n \equiv 5 \pmod{8}$.

 (2) $2^n \equiv n \pmod{125}$. This is extremely computationally intensive if we try to calculate all $2^1,2^2,\ldots,2^{100} \pmod{125}$, so we take a divide-and-conquer approach instead. In order for this expression to be true, $2^n \equiv n \pmod{5}$ is necessary; it shouldn't take too long for us to go through the $20$ possible LHS-RHS combinations, considering that $n$ must be odd. We only need to test $10$ values of $n$ and obtain $n \equiv 3 \pmod{20}$ or $n \equiv 17 \pmod{20}$.

 With this in mind we consider $2^n \equiv n \pmod{25}$. By the Generalized Fermat's Little Theorem, $2^{20} \equiv 1 \pmod{25}$, but we already have $n$ modulo $20$. Our calculation is greatly simplified. The LHS's cycle length is $20$ and the RHS's cycle length is $25$, from which their least common multiple is $100$. In this step we need to test all the numbers between $1$ to $100$ that $n \equiv 3 \pmod{20}$ or $n \equiv 17 \pmod{20}$. In the case that $n \equiv 3 \pmod{20}$, the RHS goes $3,23,43,63,83$, and we need $2^n \equiv n \equiv 2^3 \pmod{25}$; clearly $n \equiv 83 \pmod{100}$. In the case that $n \equiv 17 \pmod{20}$, by a similar argument, $n \equiv 97 \pmod{100}$. 

 In the final step, we need to calculate $2^{97}$ and $2^{83}$ modulo $125$:

 Note that $2^{97}\equiv2^{-3}$; because $8\cdot47=376\equiv1\pmod{125},$ we get $2^{97} \equiv 47\pmod{125}$. 

 Note that $2^{83}$ is $2^{-17}=2^{-16}\cdot2^{-1}$. We have
 \begin{align*} 2^{-1}&\equiv63, \\ 2^{-2}&\equiv63^2=3969\equiv-31, \\ 2^{-4}&\equiv(-31)^2=961\equiv-39, \\ 2^{-8}&\equiv1521\equiv21, \\ 2^{-16}&\equiv441, \\ 2^{-17}&\equiv63\cdot441\equiv7\cdot(-31)=-217\equiv33. \end{align*}
 This time, LHS cycle is $100$, RHS cycle is $125$, so we need to figure out $n$ modulo $500$. It should be $n \equiv 283,297 \pmod{500}$.

 Put everything together. By the second subcondition, the only candidates less than $1000$ are $283,297,783,797$. Apply the first subcondition, $n=\boxed{797}$ is the desired answer.

 





\textbf{Solution 2}

We have that $2^n + 5^n \equiv n\pmod{1000}$, or $2^n + 5^n \equiv n \pmod{8}$ and $2^n + 5^n \equiv n \pmod{125}$ by CRT. It is easy to check $n < 3$ don't work, so we have that $n \geq 3$. Then, $2^n \equiv 0 \pmod{8}$ and $5^n \equiv 0 \pmod{125}$, so we just have $5^n \equiv n \pmod{8}$ and $2^n \equiv n \pmod{125}$. Let us consider both of these congruences separately.

 First, we look at $5^n \equiv n \pmod{8}$. By Euler's Totient Theorem (ETT), we have $5^4 \equiv 1 \pmod{8}$, so $5^5 \equiv 5 \pmod{8}$. On the RHS of the congruence, the possible values of $n$ are all nonnegative integers less than $8$ and on the RHS the only possible values are $5$ and $1$. However, for $5^n$ to be $1 \pmod{8}$ we must have $n \equiv 0 \pmod{4}$, a contradiction. So, the only possible values of $n$ are when $n \equiv 5 \pmod{8} \implies n = 8k+5$.

 Now we look at $2^n \equiv n \pmod{125}$. Plugging in $n = 8k+5$, we get $2^{8k+5} \equiv 8k+5 \pmod{125} \implies 2^{8k} \cdot 32 \equiv 8k+5 \pmod{125}$. Note, for $2^n \equiv n\pmod{125}$ to be satisfied, we must have $2^n \equiv n \pmod{5}$ and $2^n \equiv n\pmod{25}$. Since $2^{8k} \equiv 1\pmod{5}$ as $8k \equiv 0\pmod{4}$, we have $2 \equiv -2k \pmod{5} \implies k = 5m-1$. Then, $n = 8(5m-1) + 5 = 40m-3$. Now, we get $2^{40m-3} \equiv 40m-3 \pmod{125}$. Using the fact that $2^n \equiv n\pmod{25}$, we get $2^{-3} \equiv 15m-3 \pmod{25}$. The inverse of $2$ modulo $25$ is obviously $13$, so $2^{-3} \equiv 13^3 \equiv 22 \pmod{25}$, so $15m \equiv 0 \pmod{25} \implies m = 5s$. Plugging in $m = 5s$, we get $n = 200s - 3$.

 Now, we are finally ready to plug $n$ into the congruence modulo $125$. Plugging in, we get $2^{200s-3} \equiv 200s - 3 \pmod{125}$. By ETT, we get $2^{100} \equiv 1 \pmod{125}$, so $2^{200s- 3} \equiv 2^{-3} \equiv 47 \pmod{125}$. Then, $200s \equiv 50 \pmod{125} \implies s \equiv 4 \pmod{5} \implies s = 5y+4$. Plugging this in, we get $n = 200(5y+4) - 3 = 1000y+797$, implying the smallest value of $n$ is simply $\boxed{797}$. 

 





\textbf{Solution 3 (Chinese Remainder Theorem and Binomial Theorem)}

We wish to find the least positive integer $n$ for which $2^n+5^n-n\equiv0\pmod{1000}.$ Rearranging gives 
 \[2^n+5^n\equiv n\pmod{1000}.\] Applying the Chinese Remainder Theorem, we get the following system of congruences:
 \begin{align*} 2^n+5^n &\equiv n \pmod{8}, \\ 2^n+5^n &\equiv n \pmod{125}. \end{align*}
It is clear that $n\geq3,$ from which we simplify to 
 \begin{align*} 5^n &\equiv n \pmod{8}, \hspace{15mm} &(1) \\ 2^n &\equiv n \pmod{125}. &(2) \end{align*}
We solve each congruence separately:

 <ol style="margin-left: 1.5em;"> <li>For $(1),$ quick inspections produce that $5^1,5^2,5^3,5^4,\ldots$ are congruent to $5,1,5,1,\ldots$ modulo $8,$ respectively. More generally, $5^n \equiv 5 \pmod{8}$ if $n$ is odd, and $5^n \equiv 1 \pmod{8}$ if $n$ is even. As $5^n$ is always odd (so is $n$), we must have $n\equiv5\pmod{8}.$ <i>\textbf{That is, $\boldsymbol{n=8r+5}$ for some nonnegative integer $\boldsymbol{r.}$}</i>

 </li> 

 \item For $(2),$ we substitute the result from $(1)$ and simplify:
 \begin{align*} 2^{8r+5}&\equiv8r+5\pmod{125} \\ \left(2^8\right)^r\cdot2^5&\equiv8r+5\pmod{125} \\ 256^r\cdot32&\equiv8r+5\pmod{125} \\ 6^r\cdot32&\equiv8r+5\pmod{125}. \end{align*}
Note that $5^3=125$ and $6=5+1,$ so we apply the Binomial Theorem to the left side:
 \begin{align*} (5+1)^r\cdot32&\equiv8r+5\pmod{125} \\ \Biggl[\binom{r}{0}5^0+\binom{r}{1}5^1+\binom{r}{2}5^2+\phantom{ }\underbrace{\binom{r}{3}5^3+\cdots+\binom{r}{r}5^r}_{0\pmod{125}}\phantom{ }\Biggr]\cdot32&\equiv8r+5\pmod{125} \\ \left[1+5r+\frac{25r(r-1)}{2}\right]\cdot32&\equiv8r+5\pmod{125} \\ 32+160r+400r(r-1)&\equiv8r+5\pmod{125} \\ 32+35r+25r(r-1)&\equiv8r+5\pmod{125} \\ 25r^2+2r+27&\equiv0\phantom{r+5}\pmod{125}. \hspace{15mm} (*) \end{align*}
Since $125\equiv0\pmod{5},$ it follows that
 \begin{align*} 25r^2+2r+27&\equiv0\pmod{5} \\ 2r+2&\equiv0\pmod{5} \\ r&\equiv4\pmod{5}. \end{align*}
<i>\textbf{That is, $\boldsymbol{r=5s+4}$ for some nonnegative integer $\boldsymbol{s.}$}</i>Substituting this back into $(*),$ we get
 \begin{align*} 25(5s+4)^2+2(5s+4)+27&\equiv0\pmod{125} \\ 625s^2+1010s+435&\equiv0\pmod{125} \\ 10s+60&\equiv0\pmod{125} \\ 10(s+6)&\equiv0\pmod{125} \\ s+6&\equiv0\pmod{25}. \end{align*}
<i>\textbf{Therefore, the least such nonnegative integer $\boldsymbol{s}$ is $\boldsymbol{19.}$}</i>

 </ol>Finally, combining the bolded results from above generates the least such positive integer $n$ at $s=19:$
 \begin{align*} n&=8r+5 \\ &=8(5s+4)+5 \\ &=40s+37 \\ &=\boxed{797}. \end{align*}
~MRENTHUSIASM (inspired by Math Jams's \textbf{2021 AIME II Discussion})

 





\textbf{Solution 4 (Minimal Computation)}


 \[5^n \equiv n \pmod{8}\]Note that $5^n \equiv 5,1,5,1,...$ and $n \equiv 0,..,7$ so $n$ is periodic every $[2,8]=8$. Easy to check that only $n\equiv 5 \pmod{8}$ satisfy. Let $n=8a+5$. Note that by binomial theorem, 
 \[2^{8a+5}=32\cdot 2^{8a} \equiv 7(1+15)^{2a} \equiv 7(1+30a)\pmod{25}\]So we have $7(1+30a) \equiv 8a+5\pmod{25} \implies 202a \equiv 23 \implies 2a \equiv 23+25 \implies a \equiv 24 \pmod{25}$.Combining $a\equiv 24\pmod{25}$ with $n \equiv 5 \pmod{8}$ gives that $n$ is in the form of $197+200k$, $n=197,397,597,797$. Note that since $\phi(125)=100$
 \[2^{197+200k} \equiv 2^{197} \equiv 2^{-3} \equiv \underbrace{47}_{\text{Extended Euclidean Algorithm}}\pmod{125}\]Easy to check that only $\boxed{797}\equiv 47\pmod{125},{1797}\equiv 47\pmod{125},{2797}\equiv 47\pmod{125},...$

 





\textbf{Solution 5 (STEP by step transform)}

1. The desired $n$ has the form  $n = m + 20l + 100p,$ where  $m, l, p$ are integers and $m < 20.$

 Really: 
 \[(2^{m+ 20} – (m + 20)) – (2^m – m) =  (2^{m+ 20} – 2^m) – 20.\] The first term is a multiple of $100$(<i>\textbf{Claim}</i>).

 We denote <i>\textbf{step}</i> an increase in $m$ by 20. At each <i>\textbf{step}</i>, the divisibility by $10$ is preserved, the number of tens is reduced by $2$.

 
 \[(2^{m+ 100} – (m + 100)) – (2^m – m) =  (2^{m+ 100} – 2^m) – 100.\] We denote <i>\textbf{STEP}</i> an increase in $m$ by $100.$ At each <i>\textbf{STEP}</i>, the first term is a multiple of $1000,$ which means that at each <i>\textbf{STEP}</i> the divisibility by $100$ is preserved, the number of hundreds decreases by $1.$

 2. If the expression $2^n + 5^n – n$ is a multiple of $1000,$ the number $n$ is odd ($2^n$ is even, $5^n$ is odd), which means that $5^n$ ends in $125.$ Therefore, the number $2^n – n$ ends in $875.$

 3. $2^3 – 3 = 8 – 3 = 5.$ The tens digit is even $(0),$ it cannot be transformed into $7$ in several steps, since at each <i>\textbf{step}</i> this digit changes by $2.$

 $17 = 20 – 3,$ so $2^{17} + 2^3$ is a multiple of $10,$ $(2^{17} – 17) + (2^3 – 3) = 2^{17} + 2^3 – 20$ is a multiple of $10.$ 
 \[2^{17} \pmod {100} = ((2^{10} \pmod {100}) \cdot (2^7 \pmod {100})) \pmod {100} = (24 \cdot 28) \pmod {100} = 72.\]
 \[(2^{17} – 17 ) \pmod {100}  = 55.\]The tens digit $5$ is odd, subtracting $2$ at each <i>\textbf{step}</i> in $4$ <i>\textbf{steps}</i> of $20$ will turn it into $7.$So 
 \[(2^{97} – 97 ) \pmod {100} = 75.\]
 \[2^{97} \pmod {1000} =  ((2^{49} \pmod {1000}) \cdot 2^7) \pmod {1000} =  672.\] 
 \[(2^{97} – 97 ) \pmod {1000} = 575.\]We transform $5$ into $8$ in $7$ <i>\textbf{STEPS}</i>, so  
 \[(2^{797} – 797 ) \pmod {1000} = 875.\]
 \[(5^{797} + 2^{797} – 797 ) \pmod {1000} = 0.\]Note, that for $n = 797 + 1000k, k$ is an integer, the expression $2^n + 5^n – n$ is a multiple of $1000.$

 <i>\textbf{Claim}</i>

 The numbers  $2^{n+2\cdot 5^k}+2^n=2^n(2^{2\cdot 5^k}+1)$ and $2^{n+4\cdot 5^k}-2^n=2^n(2^{4\cdot 5^k}-1)$ 

 are a multiple of $10^{k+1}$ for any $n > k$, where $n,k$ are an integer.

 <i>\textbf{Proof}</i>

 First, if $m \pmod{10} \equiv 4,$ then  $(m^4 – m^3 + m^2 – m + 1) \pmod{10} \equiv 6 – 4 + 6 – 4 + 1 = 5.$

 $2^{4n+2}\pmod{10} \equiv 4.$ So, if the number  $2^{4n+2}  + 1$ is a multiple of $5^k$, then $2^{20n+10} + 1$ is a multiple of $5^{k +1}.$

 Really, denote $m =  2^{4n+2},$ then,  $2^{20n+10} + 1 =  m^5 + 1 = (m + 1)\cdot (m^4 – m^3 + m^2 – m + 1).$

 The first factor is a multiple of $5^k$, the second factor is a multiple of $5.$ Their product is a multiple of  $5^{k +1}.$

 For odd $n,$ using Newton's binomial for $4^n$, we get $2^{2n} + 1$ is a multiple of $5.$ 
 \begin{align*} 2^{2n} + 1 = 4^n + 1 = (5 - 1)^n + 1 = 5^n -n\cdot 5^{n-1} +...+5n\end{align*}
If $n = 5^k,$ then $n > k,$ each term is  a multiple of $5n$.

 Therefore, the number $4^{5^k}+1$ is  a multiple of $5^{k+1}$ and  $2^{k+1} (2^{2\cdot 5^k}+1)$ is  a multiple of $10^{k+1}$.

 The difference 
 \[2^{n+4\cdot 5^k}-2^n=2^n(2^{4\cdot 5^k}-1) = 2^n(2^{2\cdot 5^k}-1)(2^{2\cdot 5^k}+1)\] is a multiple of $10^{k+1}.$

 \textbf{vladimir.shelomovskii@gmail.com, vvsss}

 





\textbf{Solution 6 (Art of Enumeration)}

Problem require us to find minimum $n$ where $2^n + 5^n - n$ is a multiple of $1000$.

 Since we already know that multiplication does not affect mod,we can simply just focus on enumerating and hope to find some pattern.

 The pattern of $5^n$ is pretty simple, it starts with 5,25, and then repeat the pattern 125,625, in other words, $5^n\equiv 125 \mod{1000}$ when n is odd, and $5^n\equiv 625 \mod{1000}$ when n is even if we ignore $n=1$ and $n=2$.

 The pattern of $2^n$, however, is quite long. In fact, it happens when $n=103$, where $2^{103} \equiv 008 \mod{1000}$. More specifically, the pattern starts with $n=3$, which means for every 100 n, $2^n \mod 1000$ will repeat in a loop: 

 $008{,}016{,}032{,}064{,}128{,}256{,}512{,}024{,}048{,}096{,}$$192{,}384{,}768{,}536{,}072{,}144{,}288{,}576{,}152{,}304{,}$$608{,}216{,}432{,}864{,}728{,}456{,}912{,}824{,}648{,}296{,}$$592{,}184{,}368{,}736{,}472{,}944{,}888{,}776{,}552{,}104{,}$$208{,}416{,}832{,}664{,}328{,}656{,}312{,}624{,}248{,}496{,}$$992{,}984{,}968{,}936{,}872{,}744{,}488{,}976{,}952{,}904{,}$$808{,}616{,}232{,}464{,}928{,}856{,}712{,}424{,}848{,}696{,}$$392{,}784{,}568{,}136{,}272{,}544{,}088{,}176{,}352{,}704{,}$$408{,}816{,}632{,}264{,}528{,}056{,}112{,}224{,}448{,}896{,}$$792{,}584{,}168{,}336{,}672{,}344{,}688{,}376{,}752{,}504$

 Moreover, we might found that for every 20 consequtive integer, $2^n \mod 100$ will repeat in a loop. Thus, by adding $5^n$ to first twenty element, which is $n=3$ to $n=22$, we can get a pattern of last two digit of $2^n + 5^n \mod{1000}$, which is 

 $33{,}41{,}57{,}89{,}53{,}$

 $81{,}37{,}49{,}73{,}21{,}$

 $17{,}09{,}93{,}61{,}97{,}$

 $69{,}13{,}01{,}77{,}29$

 By subtracting n, we can also get following sequence:

 $30{,}37{,}52{,}83{,}46{,}$

 $73{,}28{,}39{,}62{,}09{,}$

 $04{,}95{,}78{,}45{,}80{,}$

 $51{,}94{,}81{,}56{,}07$

 Noticing that only when $n=3$ and $n=17$, $2^n+5^n-n \equiv 0 \mod{10}$, so $n=3+20k$ and $17+20k$ are our possible candidates. However, since the length of our pattern is 20, we know that $n=3+20k$ is definetely not answer, which only left us $n=17+20k$.

 Now, we calculate the last three digit of $2^n+5^n-n$ at $n=17, n=37, n=57, n=77, n=97$, which is $180{,}560{,}940{,}320{,}700$.

 Since the pattern will repeat for every 100 consequtive integer, the only thing that could possibly change last three digit is -n, noticing that only when $n=97$, $2^n+5^n-n \equiv 0 \mod 100$. So the answer only could be $n=97+100k_2$, by subtracting 700, which means adding 700 to n, we can get our final answer $n=\boxed{797}$.

 Also, by this method we can ensure that $\frac{2^n+5^n-n}{1000}$ is an integer only if $n=797+1000k$, all the calculation can be done without a calculator (about 15 minutes or so), need some patience though.

 If you are wondering, $2^{103} = 10{,}141{,}204{,}801{,}825{,}835{,}211{,}973{,}625{,}643{,}008$.

 \textbf{Note} If you are wondering, $\frac{2^{797}+5^{797}-797}{1000}$ is equal to the following $555$-digit number:

 $119{,}975{,}745{,}111{,}650{,}476{,}385{,}411{,}555{,}010{,}245{,}228{,}283{,}196{,}935{,}699{,}128{,}663{,}834{,}$$986{,}755{,}809{,}416{,}868{,}468{,}175{,}224{,}221{,}677{,}450{,}271{,}793{,}186{,}899{,}399{,}171{,}810{,}240{,}$$468{,}879{,}630{,}691{,}320{,}210{,}517{,}618{,}332{,}910{,}064{,}430{,}017{,}422{,}410{,}572{,}497{,}874{,}327{,}$$116{,}662{,}273{,}049{,}144{,}209{,}532{,}072{,}513{,}248{,}821{,}826{,}125{,}454{,}909{,}131{,}367{,}242{,}561{,}$$087{,}064{,}439{,}357{,}641{,}814{,}942{,}784{,}478{,}465{,}638{,}163{,}997{,}895{,}630{,}266{,}780{,}260{,}171{,}$$070{,}180{,}244{,}384{,}687{,}160{,}174{,}154{,}618{,}815{,}254{,}719{,}975{,}476{,}350{,}789{,}415{,}969{,}908{,}$$281{,}131{,}044{,}034{,}581{,}864{,}573{,}596{,}962{,}567{,}993{,}948{,}922{,}431{,}587{,}652{,}735{,}290{,}158{,}$$293{,}666{,}986{,}616{,}531{,}419{,}688{,}663{,}485{,}548{,}648{,}995{,}509{,}247{,}893{,}796{,}528{,}146{,}109{,}$$144{,}143{,}744{,}650{,}882{,}271{,}050{,}771{,}309{,}301{,}498{,}467{,}898{,}226{,}649{,}089{,}924{,}520{,}539{,}$$820{,}344{,}035{,}886{,}481{,}987{,}499{,}589{,}845{,}734{,}228{,}834{,}491{,}187.$

 \href{https://mathproblemsolvingskills.wordpress.com/}{MathProblemSolvingSkills.com}

 Did not use any advanced method, easily understandable:

 




	\item (\textit{AIME II 2021}) Let $\Delta ABC$ be an acute triangle with circumcenter $O$ and centroid $G$. Let $X$ be the intersection of the line tangent to the circumcircle of $\Delta ABC$ at $A$ and the line perpendicular to $GO$ at $G$. Let $Y$ be the intersection of lines $XG$ and $BC$. Given that the measures of $\angle ABC, \angle BCA,$ and $\angle XOY$ are in the ratio $13 : 2 : 17,$ the degree measure of $\angle BAC$ can be written as $\frac{m}{n},$ where $m$ and $n$ are relatively prime positive integers. Find $m+n$.


 
 \textbf{Answer: }592



\textbf{Solution 1}

In this solution, all angle measures are in degrees.

 Let $M$ be the midpoint of $\overline{BC}$ so that $\overline{OM}\perp\overline{BC}$ and $A,G,M$ are collinear. Let $\angle ABC=13k,\angle BCA=2k$ and $\angle XOY=17k.$ 

 Note that:

 <ol style="margin-left: 1.5em;"> <li>Since $\angle OGX = \angle OAX = 90,$ quadrilateral $OGAX$ is cyclic by the Converse of the Inscribed Angle Theorem.It follows that $\angle OAG = \angle OXG,$ as they share the same intercepted arc $\widehat{OG}.$

 </li> 

 <li>Since $\angle OGY = \angle OMY = 90,$ quadrilateral $OGYM$ is cyclic by the supplementary opposite angles.It follows that $\angle OMG = \angle OYG,$ as they share the same intercepted arc $\widehat{OG}.$

 </li>

 </ol>Together, we conclude that $\triangle OAM \sim \triangle OXY$ by AA, so $\angle AOM = \angle XOY = 17k.$

 Next, we express $\angle BAC$ in terms of $k.$ By angle addition, we have
 \begin{align*} \angle AOM &= \angle AOB + \angle BOM \\ &= 2\angle BCA + \frac12\angle BOC \hspace{10mm} &&\text{by Inscribed Angle Theorem and Perpendicular Bisector Property} \\ &= 2\angle BCA + \angle BAC. &&\text{by Inscribed Angle Theorem} \end{align*}
Substituting back gives $17k=2(2k)+\angle BAC,$ from which $\angle BAC=13k.$

 For the sum of the interior angles of $\triangle ABC,$ we get
 \begin{align*} \angle ABC + \angle BCA + \angle BAC &= 180 \\ 13k+2k+13k&=180 \\ 28k&=180 \\ k&=\frac{45}{7}. \end{align*}
Finally, we obtain $\angle BAC=13k=\frac{585}{7},$ from which the answer is $585+7=\boxed{592}.$

 





\textbf{Solution 2}

Let $M$ be the midpoint of $BC$. Because $\angle{OAX}=\angle{OGX}=\angle{OGY}=\angle{OMY}=90^o$, $AXOG$ and $OMYG$ are cyclic, so $O$ is the center of the spiral similarity sending $AM$ to $XY$, and $\angle{XOY}=\angle{AOM}$. Because $\angle{AOM}=2\angle{BCA}+\angle{BAC}$, it's easy to get $\frac{585}{7} \implies \boxed{592}$ from here.

 





\textbf{Solution 3 (Easy and Simple)}

Firstly, let $M$ be the midpoint of $BC$. Then, $\angle OMB = 90^o$. Now, note that since $\angle OGX = \angle XAO = 90^o$, quadrilateral $AGOX$ is cyclic. Also, because $\angle OMY + \angle OGY = 180^o$, $OMYG$ is also cyclic. Now, we define some variables: let $\alpha$ be the constant such that $\angle ABC = 13\alpha, \angle ACB = 2\alpha,$ and $\angle XOY = 17\alpha$. Also, let $\beta = \angle OMG = \angle OYG$ and $\theta = \angle OXG = \angle OAG$ (due to the fact that $AGOX$ and $OMYG$ are cyclic). Then, 
 \[\angle XOY = 180 - \beta - \theta = 17\alpha \implies \beta + \theta = 180 - 17\alpha.\] Now, because $AX$ is tangent to the circumcircle at $A$, $\angle XAC = \angle CBA = 13\alpha$, and $\angle CAO = \angle OAX - \angle CAX = 90 - 13\alpha$. Finally, notice that $\angle AMB = \angle OMB - \angle OMG = 90 - \beta$. Then, 
 \[\angle BAM = 180 - \angle ABC - \angle AMB = 180 - 13\alpha - (90 - \beta) = 90 + \beta - 13\alpha.\] Thus, 
 \[\angle BAC = \angle BAM + \angle MAO + \angle OAC = 90 + \beta - 13\alpha + \theta + 90 - 13\alpha = 180 - 26\alpha + (\beta + \theta),\] and 
 \[180 = \angle BAC + 13\alpha + 2\alpha = 180 - 11\alpha + \beta + \theta \implies \beta + \theta = 11\alpha.\] However, from before, $\beta+\theta = 180 - 17 \alpha$, so $11 \alpha = 180 - 17 \alpha \implies 180 = 28 \alpha \implies \alpha = \frac{180}{28}$. To finish the problem, we simply compute 
 \[\angle BAC = 180 - 15 \alpha = 180 \cdot \left(1 - \frac{15}{28}\right) = 180 \cdot \frac{13}{28} = \frac{585}{7},\] so our final answer is $585+7=\boxed{592}$.

 





\textbf{Solution 4 (Why Isosceles)}


                    \begin{figure}[H]
                        \centering
                        \includegraphics[width=0.5\textwidth]{Solution/AIME_2021_II_14_Solution_Figure_1.png}
                    \end{figure}
                    $\angle OAX = \angle OGX = 90^\circ \implies$ quadrilateral $XAGO$ is cyclic $\implies$

 $\angle GXO =  \angle GAO,$ as they share the same intersept  $\overset{\Large\frown} {GO}.$

 $\angle OGY = \angle OMY =  90^\circ \implies$ quadrilateral $OGYM$ is cyclic $\implies$

 $\angle GYO = \angle OMG,$ as they share the same intercept $\overset{\Large\frown} {GO}.$

 In triangles $\triangle XOY$ and  $\triangle AOM,$ two pairs of angles are equal, which means that the third angles  $\angle XOY =  \angle AOM$ are also equal.

 $\angle ABC :  \angle BCA :  \angle AOM = 13 : 2 : 17,$ so  $\angle AOM =  \angle ABC + 2 \angle BCA.$

 According to the <i>\textbf{Claim}</i>,  $\triangle ABC$ is isosceles,
 \[\angle ABC : \angle BCA : \angle BAC = 13 : 2 : 13.\]
 \[\angle BAC = \frac{13} {13 + 2 + 13} \cdot 180^\circ =  \frac {585^\circ}{7} \implies  585 + 7 = \boxed{592}.\]

 
                    \begin{figure}[H]
                        \centering
                        \includegraphics[width=0.5\textwidth]{Solution/AIME_2021_II_14_Solution_Figure_2.png}
                    \end{figure}
                    <i>\textbf{Claim}</i>

 Let $\triangle ABC$ be an acute triangle with circumcenter $O.$

 Let $M$ be the midpoint of $BC$ so $MO\perp BC.$ 

 If $\angle AOM = 2\angle ACB + \angle ABC,$ then $AC = BC.$

 We define $\angle AOM$ as the sum of $\angle AOB + \angle BOM,$ this angle can be greater than $180^\circ.$

 <i>\textbf{Proof}</i>

 $\angle BAC = \angle BOM$  as they share the same intercept $\overset{\Large\frown} {BC}$ (an inscribed angle and half of central angle). 

 $\angle AOB = 2\angle ACB$  as they share the same intercept $\overset{\Large\frown} {AB}.$

 
 \[\angle AOM = \angle AOB +  \angle BOM = 2 \angle ACB +  \angle CAB.\]

 If  $\angle AOM = 2 \angle ACB +  \angle ABC,$ then  $\angle ABC =  \angle CAB, AC = BC.$

 \textbf{vladimir.shelomovskii@gmail.com, vvsss}

 





\textbf{Solution 5}

Extend $XA$ and meet line $CB$ at $P$. Extend $AG$ to meet $BC$ at $F$. Since $AF$ is the median from $A$ to $BC$, $A,G,F$ are collinear. Furthermore, $OF$ is perpendicular to $BC$

 Draw the circumcircle of $\triangle{XPY}$, as $OA\bot XP, OG\bot XY, OF\bot PY$, $A,G,F$ are collinear, $O$ lies on the circumcircle of $\triangle{XPY}$ as $AGF$ is the Simson line of $\triangle{XPY}$. Thus, $\angle{P}=180-17x, \angle{PAB}=\angle{C}=2x, 180-15x=13x, x=\frac{45}{7}$, the answer is $180-15\cdot \frac{45}{7}=\frac{585}{7}$ which is $\boxed{592}$.

 




	\item (\textit{AIME II 2021}) Let $f(n)$ and $g(n)$ be functions satisfying

 \[f(n) =  \begin{cases} \sqrt{n} & \text{ if } \sqrt{n} \text{ is an integer}\\ 1 + f(n+1) & \text{ otherwise} \end{cases}\]
and

 \[g(n) = \begin{cases}\sqrt{n} & \text{ if } \sqrt{n} \text{ is an integer}\\ 2 + g(n+2) & \text{ otherwise} \end{cases}\]
for positive integers $n$. Find the least positive integer $n$ such that $\frac{f(n)}{g(n)} = \frac{4}{7}$.


 
 \textbf{Answer: }258



\textbf{Solution 1}

Consider what happens when we try to calculate $f(n)$ where n is not a square. If $k^2<n<(k+1)^2$ for (positive) integer k, recursively calculating the value of the function gives us $f(n)=(k+1)^2-n+f((k+1)^2)=k^2+3k+2-n$. Note that this formula also returns the correct value when $n=(k+1)^2$, but not when $n=k^2$. Thus $f(n)=k^2+3k+2-n$ for 

 




\end{enumerate}